\chapter{Reumatoidná artritída}

1 Reumatoidná artritída Reumatoidná artritída (RA) predstavuje závažné, chronické zápalové ochorenie au- toimunitnej etiológie, ktoré sa vyznacuje systémovým charakterom a progresívnym priebe- hom. Jablkom sváru v patogenéze tohto ochorenia je imunitný systém, ktorý omylom iden- tifikuje vlastné tkanivá ako cudzorodé, co vedie k deštrukcii, bolesti a strate funkcie. Hoci primárnymciel’ompatologickéhoprocesusúperiférnesynoviálneklby,najmädrobnésklbe- nia rúk a nôh, RA sa v plnej miere prejavuje ako systémová choroba postihujúca aj mimo- klbové štruktúry a vnútorné orgány. Tento aspekt ochorenia významne prispieva k zvýšenej morbiditeaskráteniudlžkyživotapostihnutýchjedincov.Ochoreniesamôžemanifestovat’v ktoromkol’vekveku,avšakdemografickékrivkyjasnepoukazujúnavrcholincidencievpro- duktívnomveku,typickymedzi30.aža50.rokomživota,cománesmiernycelospolocenský a ekonomický dopad. Etiopatogenéza ochorenia je multifaktoriálna a hoci nie je do posled- ného detailu objasnená, predpokladá sa zložitá interakcia genetickej predispozície hostitel’a a spúšt’acích faktorov vonkajšieho prostredia. Priebeh choroby je vysoko individuálny -- od miernych foriem s pomalou progresiou až po agresívne formy vedúce k rýchlej invalidizácii (Rovenský,2001). Anatomickým substrátom, kde sa odohráva hlavná dráma autoimunitnej reakcie, je synoviálna membrána. Za normálnych okolností tenká vrstva buniek sa pri RA mení na hy- perplastické,zapálenétkanivoprestúpenébunkamiimunitnéhosystému.Tátozmenenásyno- via, oznacovaná ako panus, sa správa ako lokálny tumor, ktorý invazívne prerastá do okolia, nahlodávaklbovúchrupavkuanásledneajsubchondrálnukost’.Výsledkomtohtoprocesuje nielen bolest’ a opuch, ale predovšetkým nezvratná deštrukcia klbovej architektúry a strata hybnosti(Rovenský,2001;Shahetal.,2024).


\section{Funkcnáanatómiaklbovéhosystému}

Pochopenie anatomických a biomechanických princípov klbového aparátu je nevy- hnutným predpokladom pre správnu interpretáciu patologických zmien pri reumatoidnej ar- tritíde.Klbyl’udskéhotelamožnorozdelit’podl’arôznychkritérií,pricomzpohl’adureuma- tológie má primárny význam delenie podl’a prítomnosti klbovej dutiny na spojenia nepravé (synartrózy)apravéklby(diartrózy).Právepravéklbyobsahujúcesynoviálnumembránusú tercomautoimunitnéhoútokupriRA(Dylevský,2009).


\subsection{Stavbasynoviálnehoklbu}

Synoviálny klb predstavuje zložitú anatomickú jednotku, ktorej základnými kompo- nentmi sú klbové povrchy pokryté hyalínnou chrupavkou, klbové puzdro s jeho vrstvami a klbová dutina vyplnená synoviálnou tekutinou. Každá z týchto štruktúr zohráva špecifickú úlohu v biomechanike pohybu a súcasne predstavuje potenciálny substrát pre patologické 2 zmeny(Dylevský,2009). Klbová chrupavka pokrýva artikulacné plochy kostí a zabezpecuje hladký pohyb s minimálnymtrením.JetvorenáprevažneextracelulárnoumatrixbohatounakolagéntypuIIa proteoglykány,ktoréviažuvoduaposkytujúchrupavkejejtypickéviskoelastickévlastnosti. Chondrocyty, jediné bunky prítomné v chrupavke, udržiavajú rovnováhu medzi syntézou a degradácioumatrix.PriRAdochádzaknarušeniutejtorovnováhyvprospechdegradácie,co vediekprogresívnemuúbytkuchrupavkovéhotkaniva. Klbovépuzdropozostávazdvochvrstiev.Vonkajšiafibróznavrstva(membranafib- rosa) je tvorená hustým väzivom bohatým na kolagénové vlákna typu I, ktoré zabezpecujú mechanickú stabilitu klbu. Vnútorná synoviálna vrstva (membrana synovialis) je pri RA miestom najintenzívnejších zápalových zmien. Za fyziologických podmienok je synoviálna membránatenká,tvorenájednouažtromivrstvamibuniek--synoviocytov(Kolár,2009).


\subsection{Synoviálnamembránaajejúloha}

Synoviálna membrána je metabolicky aktívne tkanivo, ktoré produkuje synoviálnu tekutinu -- viskózny ultrafiltrátt krvnej plazmy obohatený o kyselinu hyalurónovú. Táto te- kutinaplníniekol’kokl’úcovýchfunkcií:lubrikáciuklbovýchpovrchov,výživuavaskulárnej chrupavkydifúziouživínaodstranovaniemetabolickýchproduktovzklbovejdutiny. V synoviálnej membráne rozlišujeme dva hlavné typy buniek. Synoviocyty typu A majúmakrofágovýpôvodazabezpecujúfagocytózudebrisaimunitnýdohl’ad.Synoviocyty typu B sú fibroblastového charakteru a produkujú komponenty synoviálnej tekutiny, najmä kyselinu hyalurónovú a lubricín. Pri RA dochádza k masívnej infiltrácii synoviálnej mem- brányimunitnýmibunkamiatransformáciirezidentnýchsynoviocytovdoagresívnehofeno- typu(Smith,2015).


\subsection{AnatomicképredilekcnémiestapriRA}

Reumatoidná artritída vykazuje charakteristickú predilekciu pre urcité klbové sku- piny, co má diagnostický význam. Najcastejšie postihnuté sú metakarpofalangeálne (MCP) klby,proximálneinterfalangeálne(PIP)klbyrúk,zápästnéklbyametatarzofalangeálne(MTP) klby nôh. Distálne interfalangeálne (DIP) klby sú typicky ušetrené, co odlišuje RA od oste- oartrózyapsoriatickejartritídy(Olejárová,2008). Táto selektívna distribúcia postihnutia nie je náhodná. Drobné klby prstov obsahujú pomerne vel’kú synoviálnu membránu v pomere k vel’kosti klbu, co poskytuje väcšiu plo- chupreautoimunitnúreakciu. Navyšetietoklbysúvystavenéopakovanému mechanickému stresu pri bežných denných cinnostiach, co môže potencovat’ zápalový proces (K. Pavelka; Rovenskýetal.,2003). 3


\section{Epidemiológia}

Výskyt a rozšírenie RA sa líši podl’a diagnostických schém, pricom incidencia je vyššia u žien než u mužov, a to v 70 percentách prípadov. Prevalencia sa pohybuje okolo pol percenta až jedného percenta populácie, líši sa avšak podl’a regiónu a etnických skupín. Zatial’ co v Amerike je prevalencia vyššia, v Afrike je práve naopak nižšia. RA je charakte- ristickávýskytomvoveku20až55rokov;snovýmištúdiamisapotvrdilo,ževznikáužajv neskoršomštádiuveku(Rovenský,2001;Rovenský,1998). Podl’a odhadov v roku 2020 trpelo RA celosvetovo približne 17,6 milióna obyvate- l’ov, po zohl’adnení vekovej štruktúry populácie sa odhaduje, že celosvetovo postihuje re- umatoidnáartritídapribližne0,21%l’udí.Vporovnanísrokom1990došlok14,1%nárastu poctuprípadovRA.Predpokladása,žedo25rokovpocetpacientovsRAstúpnena31,7mi- lióna. Epidemiologické štúdie v Cíne uvádzajú okolo 5 miliónov pacientov, co predstavuje prevalenciu0,42%populácie(Gaoetal.,2020;Rovenský,2001;Rovenský,1998). U starších pacientov zaznamenávame odlišné epidemiologické charakteristiky. Ge- rontorevmatológia poukazuje na špecifiká RA v geriatrickej populácii, kde ochorenie casto nadobúda atypický klinický priebeh. Séropozitivita býva nižšia a diferenciálna diagnostika vocipolymyalgiireumatikeaosteoartrózejenárocnejšia(Rovenský,2014).


\section{Etiológiaapatofyziológia}

Súcasná veda nazerá na vznik reumatoidnej artritídy ako na viacstupnový proces, v ktorom dochádza k prelomeniu imunologickej tolerancie. Na zaciatku stojí geneticky pre- disponovaný jedinec (nosic urcitých alel HLA-DRB1), u ktorého environmentálny spúš- t’ac --napríkladfajcenie,chronickáparodontitída(baktériaPorphyromonasgingivalis)alebo zmenyvcrevnommikrobióme--naštartujeprodukciuautoprotilátok.Tentopredklinickýstav môžetrvat’rokypredtým,nežsaobjavíprvásynovitída.Kl’úcovýmibiomarkerovýmihrácmi v tomto procese sú reumatoidný faktor (RF) a protilátky proti cyklickému citrulínovanému peptidu (anti-CCP), ktorých prítomnost’ casto predchádza klinickým symptómom (Shah et al.,2024;Schereretal.,2020).


\subsection{Autoimunitnémechanizmy}

Autoimunita predstavuje stav, pri ktorom imunitný systém stráca schopnost’ rozli- šovat’ medzi vlastnými a cudzorodými antigénmi. Pri fyziologických podmienkach existujú viaceré kontrolné mechanizmy -- centrálna a periférna tolerancia -- zabranujúce útoku na vlastné tkanivá. U pacientov s RA dochádza k zlyhaniu týchto mechanizmov na viacerých úrovniach(Buc,2016). Genetická predispozícia zohráva kl’úcovú úlohu v rozvoji autoimunitného procesu. Asociácia s hlavným histokompatibilným komplexom (MHC), konkrétne s urcitými ale- lami HLA-DRB1 obsahujúcimi tzv. zdiel’aný epitop (shared epitope), je jedným z najsil- nejších genetických rizikových faktorov. Tieto molekuly HLA prezentujú autoantigény T- 4 lymfocytom,címiniciujúpatologickúimunitnúodpoved’(Taketal.,2011). Proces citrulinácie predstavuje ústredný mechanizmus v patogenéze RA. Ide o po- sttranslacnú modifikáciu proteínov, pri ktorej enzým peptidylarginín deimináza (PAD) kon- vertuje arginín na citrulín. Citrulínované proteíny sa stávajú neoantigénmi, voci ktorým or- ganizmus vytvára protilátky (ACPA). Tieto autoprotilátky sa môžu objavit’ roky pred mani- festáciouklinickýchpríznakovaichprítomnost’jesilnýmprediktoromerozívnehopriebehu ochorenia(Choy,2012).


\subsection{Patofyziologickákaskáda}

V patofyziologickej kaskáde dominuje dysregulácia imunitného systému. Dochádza k patologickej aktivácii T-lymfocytov (najmä Th1 a Th17 subpopulácií), ktoré migrujú do klbovej výstelky. Tu prostredníctvom medzibunkovej komunikácie stimulujú makrofágy, B- lymfocyty a fibroblasty k nadprodukcii prozápalových mediátorov. Vzniká tzv. cytokínová búrka, v ktorej dominujú tumor nekrotizujúci faktor alfa (TNF-α), interleukín-1 (IL-1) a interleukín-6 (IL-6). Tieto molekuly sú zodpovedné nielen za lokálny zápal a bolest’, ale aj za systémové prejavy, ako je únava, horúcka ci tvorba reaktantov akútnej fázy v peceni (CRP)(Gaoetal.,2020).


\subsection{Úlohasynoviálnychfibroblastov}

Synoviálne fibroblasty,oznacované ajako fibroblastom podobné synoviocyty(FLS), zohrávajúcentrálnuúlohuvpatogenézeRA.Narozdielodimunitnýchbuniek,ktorédoklbu infiltrujú zvonka, FLS sú rezidentné bunky synoviálnej membrány, ktoré pri RA podliehajú fenotypovej transformácii. Nadobúdajú agresívne, nádoru podobné vlastnosti charakterizo- vanézvýšenouproliferáciou,zníženouapoptózouaschopnost’ouinvadovat’dochrupavkya kosti(Smith,2015). Aktivované FLS produkujú široké spektrum deštruktívnych enzýmov, predovšetkým matrix metaloproteinázy (MMP-1, MMP-3, MMP-13), ktoré degradujú kolagén a d’alšie komponenty extracelulárnej matrix. Súcasne sekretujú chemokíny prit’ahujúce d’alšie zápa- lovébunkyarastovéfaktorypodporujúceangiogenézuvzapálenomtkanive.Tentoautokriný aparakrinýkomunikacnýsystémvytvárasebaposilnujúcicykluschronickéhozápalu(Smith, 2015;Taketal.,2011).


\subsection{Tvorbapanusuadeštrukciaklbu}

Chronická stimulácia vedie k transformácii synoviocytov, ktoré sa zacínajú nekon- trolovane množit’ a správat’ invazívne. Vytvára sa hypervaskularizované granulacné tka- nivo -- panus. Bunky panusu produkujú metaloproteinázy a iné enzýmy, ktoré doslova "trá- via"klbovú chrupavku. Súcasne dochádza k aktivácii osteoklastov, buniek odbúravajúcich kost’, co vedie k vzniku typických kostných erózií. Tento deštruktívny proces je nezvratný; raz znicenú chrupavku a kost’ organizmus nedokáže plnohodnotne regenerovat’,co podciar- 5 kujenutnost’vcasnejaagresívnejliecby(Rovenský,2001;Isaacsetal.,2002).


\section{Klinickýobraz}

Klinická manifestácia RA je povestná svojou heterogenitou. U väcšiny pacientov (približne 70\%) sa ochorenie vkráda do života plazivo, nenápadne, v priebehu niekol’kých týždnovcimesiacov.Prvýmipríznakminemusiabyt’bolestiklbov,alenešpecifické,prodro- málne symptómy pripomínajúce virózu: nmerná únava, subfebrílie, nechutenstvo a mierny úbytok hmotnosti. Až následne sa pridružuje typická artikulárna symptomatológia. Domi- nantným znakom je ranná stuhnutost’ klbov. Na rozdiel od krátkodobej stuhnutosti pri ar- tróze, pri RA tento stav trvá dlhšie než hodinu, casto celé predpoludnie, a odznieva až po rozhýbaní. Stuhnutost’ je odrazom nahromadenia zápalového edému v klbe pocas nocného pokoja(Nemec,2023). Artritídasaklinickyprejavujeklasickoupäticoupríznakovzápalu:bolest’ou(dolor), opuchom (tumor), zvýšenou teplotou klbu (calor) a poruchou funkcie (functio laesa); zacer- venanie (rubor) býva pri RA menej casté. Postihnutie je typicky polyartikulárne (zasahuje viac ako 3 klbové oblasti) a symetrické (postihuje rovnaké klby na oboch stranách tela). Predilekcnýmilokalitami súdrobnéklbyruky --metakarpofalangeálne(MCP) aproximálne interfalangeálne (PIP) klby, ako aj zápästia a metatarzofalangeálne (MTP) klby nohy. Dis- tálne interfalangeálne klby (DIP) bývajú typicky ušetrené, co je dôležitý diferenciacný znak oprotiosteoartrózeapsoriatickejartritíde(Nemec,2023;Rovenský,2001). U menšej casti pacientov môže mat’ choroba akútny, búrlivý nástup, ktorý pacienta doslova pripúta na lôžko v priebehu niekol’kých dní. Ešte zriedkavejší je palindromatický reumatizmus, charakterizovaný krátkymi, opakujúcimi sa záchvatmi artritídy, ktoré nezane- chávajú trvalé následky, no môžu prejst’ do klasickej RA. Zvláštnu pozornost’ si vyžaduje RA u starších l’udí, ktorá sa casto prezentuje asymetrickým postihnutím vel’kých klbov, najmäramennýchpletencov(tzv."shoulder-likeönset),comôževiest’kzámenespolymyal- gioureumaticou(Nemec,2023).


\subsection{Klbovédeformityaštrukturálnezmeny}

Dlhodobo neostatocne kontrolovaný zápal vedie k deštrukcii väzivového aparátu a klbových puzdier, co má za následok vznik typických deformít. Tieto zmeny nie sú len koz- metickým defektom, ale vedú k t’ažkej funkcnej poruche úchopu a lokomocnej schopnosti (Olejárová;kol.,2016). Rukaazápästie:ArtikulárnykomplexrukyjezrkadlomRA.Synovitídazápästiave- diekdeštrukciikarpálnychkosticiekaväzov,cocastovyústidovolárnejsubluxáciezápästia a instability. Castou komplikáciou je kompresia *nervus medianus* -- syndróm karpálneho tunela. Na prstoch pozorujeme charakteristickú ulnárnu deviáciu (vybocenie prstov smerom klakt’ovejkosti)vMCPklboch.Dalšímitypickýmideformitamisú: • Deformitagombíkovejdierky(Boutonnière):Jecharakterizovanáfixovanouflexiou 6 vPIPklbeahyperextenziouvDIPklbe.Vznikápretrhnutímcentrálnehopruhuexten- zorovejšl’achy. • Deformita labutej šije (Swan neck): Opacná situácia -- hyperextenzia v PIP klbe a flexia v DIP klbe, spôsobená spazmom vnútorných svalov ruky a laxicitou volárnej platnicky. Noha: Postihnutie nôh je casté a bolestivé. Dochádza k poklesu priecnej klenby, vznikuhalluxvalgus(vybocenýpalec)akladivkovýchprstov(digitimallei).Podhlavickami metatarzov sa tvoria bolestivé otlaky, ktoré pacientovi st’ažujú chôdzu a vyžadujú ortope- dickúobuv. Vel’ké nosné klby: Koleno je postihnuté casto, s tvorbou rozsiahlych výpotkov. Hy- pertrofickásynoviaalebotekutinasamôžepretlacit’dopodkolennejjamkyavytvorit’Bake- rovucystu,ktorejprasknutieimitujetrombózužíllýtka.Bedrovýklbbývapostihnutýmenej casto,alekoxitídamátendenciukrýchlejprogresiiavyžadujecastototálnuendoprotézu. Chrbtica:Kritickouoblast’oujekrcnáchrbtica.Zápalvokolíatlantoaxiálnehosklbe- nia(C1-C2)môževiest’kuvol’neniuväzovainstabilite.Tentostavježivotohrozujúci,pre- tože pri neopatrnom pohybe môže dôjst’ k útlaku miechy. Každý pacient s RA podstupujúci operáciuvcelkovejanestéziimusímat’RTGkrcnejchrbticenavylúcenietejtokomplikácie. Obr. 1.1: Schematické znázornenie typických deformít ruky pri pokrocilej reumatoidnej ar- tritíde.ZobrazenájeulnárnadeviáciavMCPklbochadeformityprstov(Netter,2014) 7


\subsection{Mimoklbové(extraartikulárne)prejavy}

Reumatoidnáartritídajesystémovéochorenieparexcellence.Prítomnost’mimoklbo- výchprejavovsignalizujet’ažšiuformuochoreniaahoršiuprognózu(Rovenský,2001). Reumatoidné uzly: Sú najcastejším extraartikulárnym prejavom, vyskytujúcim sa u 20-30\% séropozitívnych pacientov. Sú to tuhé, nebolestivé podkožné útvary, ktoré sa tvoria najmä v miestach mechanického tlaku -- nad lakt’ovým výbežkom (olecranon), na chrbte rukycinaachilovke. Ocné postihnutie: Najbežnejšia je keratokonjunktivitída sicca (syndróm suchého oka) ako súcast’ sekundárneho Sjögrenovho syndrómu. Závažnejšou komplikáciou je sk- leritídaaleboepiskleritída,ktorásaprejavujebolestivýmzacervenanímokaamôževiest’až kpoškodeniuzraku. Pl’úcne postihnutie: Pl’úca sú castým tercom. Môže íst’ o pleuritídu (zápal pohrud- nice) s výpotkom, alebo o závažnejšiu intersticiálnu pl’úcnu fibrózu, ktorá postupne znižuje dýchaciukapacitupl’úc. Vaskulitída: Reumatoidná vaskulitída je zápal ciev, ktorý môže viest’ k tvorbe kož- ných vredov, nekróz na prstoch alebo k postihnutiu nervov (mononeuritis multiplex). Je to vážnakomplikáciavyžadujúcaagresívnuliecbu.


\section{Diagnostika}

Modernáreumatológiakladiedôraznavcasnúdiagnostiku,tzv.öknopríležitosti"(window ofopportunity).Ciel’omjezachytit’ochorenievjehopociatkoch,kedyjeterapeutickáinter- vencia najúcinnejšia a dokáže zabránit’ trvalému poškodeniu klbov. Diagnostický proces je komplexnýavychádzazosyntézyklinickéhoobrazu,laboratórnychvýsledkovazobrazova- cíchmetód.ZlatýmštandardomsúvsúcasnostiklasifikacnékritériáACR/EULARzroku 2010,ktorébolivyvinutépráveprezáchytvcasnýchštádiíartritídy(Aletahaetal.,2010). Bodovací systém kritérií hodnotí štyri domény, pricom na potvrdenie diagnózy "jed- noznacnejRA"jepotrebnýsúcet6aviacbodovzmaximálnych10: 1. Postihnutie klbov (0 -- 5 bodov): Dôraz sa kladie na pocet a typ postihnutých klbov. Najvyššie bodové ohodnotenie (5 bodov) získava pacient s postihnutím viac ako 10 klbov,pricomaspon jedenznichmusíbyt’malýklb. 2. Sérológia (0 -- 3 body): Prítomnost’ reumatoidných faktorov (RF) a protilátok proti citrulínovaným peptidom (ACPA/anti-CCP). Vysoké titre týchto protilátok sú vysoko špecificképreRAaprognostickynepriaznivé. 3. Reaktanty akútnej fázy (0 -- 1 bod): Laboratórny dôkaz systémového zápalu -- zvý- šenásedimentáciaerytrocytov(FW)aleboC-reaktívnyproteín(CRP). 4. Trvaniesymptómov(0--1bod):Pretrvávaniet’ažkostídlhšieako6týždnovodlišuje chronickúartritíduodprechodnýchvírusovýchartralgií. 8 Okrem klasického röntgenového vyšetrenia (RTG), ktoré v skorom štádiu môže byt’ negatívne,sadnesdopoprediadostávajúsenzitívnejšiemetódy.Ultrasonografia(USG)klbov s Power Doppler módom dokáže zobrazit’ synoviálny zápal a prekrvenie skôr, než dôjde k poškodeniukosti.Magnetickárezonancia(MRI)jenajcitlivejšoumetódounadetekciukost- néhoedému,ktorýjepredzvest’oubudúcicherózií(Smolenetal.,2018).


\section{Komplikáciereumatoidnejartritídy}

Jenevyhnutnévnímat’RAnielenakoizolovanúartropatiu,aleakokomplexnúsysté- movú chorobu, ktorej zápalový potenciál indukuje rad závažných komorbidít a komplikácií, majúcich zásadný vplyv na morbiditu a predcasnú mortalitu postihnutých jedincov. Tieto nežiaduce stavy môžu pramenit’ priamo z aktivity základného ochorenia alebo môžu byt’ iatrogénnymnásledkomdlhodobejimunosupresívnejliecby(Smolenetal.,2018).


\subsection{Kardiovaskulárnekomplikácie}

Epidemiologické štúdie konzistentne potvrdzujú, že populácia s RA celí markantne vyššiemurizikukardiovaskulárnychpríhod(približne1,5až2-násobne)vporovnanísozdra- vou populáciou. Práve kardiovaskulárne zlyhanie je hlavnou prícinou skrátenia dlžky života u týchto pacientov. Pretrvávajúci systémový zápal poškodzuje endotelovú výstelku ciev a akceleruje aterosklerotické procesy, co vyúst’uje do zvýšenej incidencie infarktu myokardu, cievnychmozgovýchpríhodakongestívnehosrdcovéhozlyhania(Smolenetal.,2018). Mechanizmykardiovaskulárnehopoškodenia Patofyziologický vzt’ah medzi reumatoidnou artritídou a kardiovaskulárnym ocho- rením je komplexný a zahrna viaceré mechanizmy. Chronický systémový zápal prispieva k endotelovejdysfunkcii--pociatocnémuštádiuaterosklerózy.Prozápalovécytokíny(TNF-α, IL-6) aktivujú endotelové bunky, cím podporujú adhéziu monocytov a ich transmigráciu do cievnej steny. Výsledkom je formácia aterómových plátov, ktoré sú u pacientov s RA casto menejstabilnéanáchylnejšienaruptúru(Freedmanetal.,2014). Rizikový profil je navyše zhoršený prítomnost’ou konvencných faktorov, ako sú hy- pertenziaadyslipidémia,ktorémôžubyt’potencovanéterapioukortikoidmianesteroidnými antireumatikami(NSA).Zklinickéhohl’adiskajevýznamnáajperikarditída,ktoráhocicasto prebieha asymptomaticky, môže byt’ detegovaná echokardiograficky. Striktný manažment krvnéhotlakualipidovjepretoimperatívom(Smolenetal.,2018;Freedmanetal.,2014). Preventívnestratégie PrevenciakardiovaskulárnychpríhodupacientovsRAvyžadujekomplexnýprístup. Na prvom mieste je efektívna kontrola zápalovej aktivity základného ochorenia -- pacienti v remisii majú významne nižšie kardiovaskulárne riziko. Dalej je potrebný striktný manaž- menttradicnýchrizikovýchfaktorovvrátaneliecbyhypertenzie,dyslipidémieadiabetu.Od- 9 porúca sa používat’ modifikované skórovacie systémy (SCORE) s korekcným faktorom 1,5 prepacientovsRA.Významnájeajúlohapravidelnejfyzickejaktivity,ktoránielenzlepšuje kardiovaskulárnukondíciu,alemáajprotizápalovéúcinky(Freedmanetal.,2014).


\subsection{Osteoporózaafraktúry}

Demineralizáciakostnéhotkaniva,osteoporóza,predstavujejednuznajcastejšíchex- traartikulárnych komplikácií s komplexnou etiológiou. V patogenéze rozlišujeme lokálnu periartikulárnu osteoporózu, ktorá je priamym následkom pôsobenia prozápalových cyto- kínov v okolí postihnutého klbu, a generalizovanú systémovú osteoporózu. Na jej vzniku sa podiel’a synergický efekt chronického zápalu, imobilizácie pacienta (zníženie mechanic- kej zát’aže kosti pre bolest’) a najmä dlhodobej kortikoterapie, ktorá inhibuje osteoblastickú aktivitu a redukuje intestinálnu resorpciu vápnika. Dôsledkom je zvýšená fragilita skeletu a vysokérizikopatologickýchfraktúr,predovšetkýmstavcovaproximálnehofemuru,covedie kd’alšiemuzhoršeniufunkcnéhostavu(KarelPavelkaetal.,2018).


\subsection{Infekcnékomplikácie}

Dysfunkcia imunitného dohl’adu, ktorá je inherentnou súcast’ou patofyziológie RA, v kombinácii s úcinkami potentných imunosupresív (vrátane biologík), robí pacientov zra- nitel’nými voci infekcným agens. Infekcie respiracného traktu (pneumónie), urogenitálneho systému a kože majú u týchto pacientov tendenciu k t’ažšiemu priebehu a castejším kom- plikáciám. Pri liecbe inhibítormi TNF-α existuje špecifické riziko reaktivácie latentnej tu- berkulózyalebochronickejhepatitídyB,covyžadujeprísnyskríningpredzahájenímliecby. Zvýšená je aj incidencia herpes zoster. Vakcinácia proti chrípke a pneumokokom je preto kl’úcovoupreventívnoustratégiou(Macejová,2019).


\subsection{Pl’úcneahematologickékomplikácie}

Pl’úcne prejavy sú castou viscerálnou manifestáciou RA. Môžu mat’ formu intersti- ciálnejpl’úcnejchoroby(ILD),ktoráprogredujedonezvratnejfibrózy,pleurálnychvýpotkov aleboreumatoidnýchuzlovvpl’úcnomparenchýme(vkombináciispneumokoniózouznáme ako Caplanov syndróm). Klinickým korulátom je progredujúca dyspnoe a suchý, dráždivý kašel’. V krvnom obraze dominuje normocytárna anémia chronických chorôb, mechaniz- momktorejjehepcidínomsprostredkovanáblokádametabolizmuželeza.Zriedkavou,noži- vot ohrozujúcou entitou je Feltyho syndróm (triáda: RA, splenomegália, granulocytopénia), ktorýextrémnezvyšujenáchylnost’nafulminantnébakteriálneinfekcie(Rovenský,2001). Intersticiálneochoreniepl’úcpriRA Intersticiálne ochorenie pl’úc asociované s reumatoidnou artritídou (RA-ILD) pred- stavujezávažnúextraartikulárnumanifestáciu,ktorávýznamnezhoršujeprognózupacientov. Prevalencia subklinického postihnutia pl’úc detegovaného pomocou vysokorozlišovacieho 10 CT (HRCT) dosahuje až 60 percent, pricom klinicky manifestná forma sa vyskytuje u 10 až 30percentchorých.Rizikovéfaktoryzahrnajúmužsképohlavie,vyššívek,fajcenie,vysoké titreACPAprotilátokaprítomnost’reumatoidnýchuzlov(Kaduraetal.,2021). PatogenézaRA-ILDzahrnakomplexnúinterakciugenetickýchfaktorov,autoimunit- ných mechanizmov a environmentálnych vplyvov. Cytokíny produkované v rámci systémo- vého zápalu, najmä TNF-alfa, IL-6 a transformujúci rastový faktor beta (TGF-β), stimulujú proliferáciu fibroblastov a nadmernú depozíciu kolagénu v pl’úcnom interstíciu. Histopato- logicky prevláda obraz obvyklej intersticiálnej pneumónie (UIP) alebo nešpecifickej inters- ticiálnejpneumónie(NSIP)(Kaduraetal.,2021;Klener,2011). Manažment RA-ILD vyžaduje multidisciplinárny prístup zahrnajúci reumatológa a pneumológa. Terapeutické možnosti sú limitované; niektoré DMARDs (napr. metotrexát) môžu potenciálne zhoršit’ pl’úcne postihnutie, zatial’ co iné (rituximab, abatacept) sa javia ako bezpecnejšie. V pokrocilých štádiách prichádza do úvahy antifibrotická liecba pirfeni- dónomalebonintedanibom(Kaduraetal.,2021).


\subsection{MetabolicképoruchypriRA}

Vzt’ahmedzireumatoidnouartritídouametabolickýmiporuchamijebidirecionálnya komplexný.ChronickýzápalpriRAnegatívneovplyvnujemetabolicképarametre,zatial’co metabolický syndróm zhoršuje zápalovú aktivitu a komplikuje liecbu základného ochorenia (Kašperováetal.,2021). Inzulínovárezistenciaadiabetesmellitus Pacienti s RA majú zvýšené riziko rozvoja diabetes mellitus 2. typu v porovnaní s bežnou populáciou. Prozápalové cytokíny, predovšetkým TNF-α a IL-6, interferujú so sig- nálnou dráhou inzulínu na úrovni inzulínového receptorového substrátu (IRS-1), co vedie k rozvoju inzulínovej rezistencie. Situáciu d’alej komplikuje dlhodobá kortikoterapia, ktorá priamozvyšujeglykémiuapodporujeviscerálnuobezitu(Kašperováetal.,2021). Dyslipidémiaaateroskleróza Lipidový profil u pacientov s aktívnou RA vykazuje charakteristické zmeny ozna- cované ako "lipidový paradox". Napriek nízkym hodnotám celkového cholesterolu a LDL- cholesterolu (ktoré sú konzumované zápalovým procesom) je kardiovaskulárne riziko zvý- šené.Prícinousúkvalitatívnezmenylipoproteínov--oxidáciaLDLcastícazníženáantiatero- génna aktivita HDL cholesterolu. Efektívna kontrola zápalovej aktivity pomocou DMARDs vedie k paradoxnému zvýšeniu hladín lipidov, co však reflektuje zlepšenie metabolického stavu(Freedmanetal.,2014;Kašperováetal.,2021). 11 Reumatickákachexia Osobitným metabolickým fenoménom pri RA je reumatická kachexia, charakteri- zovaná stratou svalovej hmoty pri zachovanej alebo dokonca zvýšenej tukovej hmote. Na rozdielodklasickejkachexieprimalígnychochoreniach,telesnáhmotnost’môžezostat’ne- zmenená, co maskuje závažný úbytok svalstva. Patofyziologicky sa uplatnuje katabolický efekt prozápalových cytokínov, znížená fyzická aktivita pre bolest’ a artritídu a niekedy aj nedostatocný príjem bielkovín. Dôsledkom je svalová slabost’, zhoršená funkcná kapacita a zvýšenérizikopádov(Klener,2011). Dopad reumatoidnej artritídy na kvalitu života (QoL) je devastacný a multodimen- zionálny, zasahujúci d’aleko za hranice fyzickej disability. WHO definuje zdravie ako stav kompletnej bio-psycho-sociálnej pohody, co je ideál, ktorý je u pacientov s RA casto nedo- siahnutel’ný(Krchnaváetal.,2018).


\subsection{Fyzickáafunkcnádoména}

Kl’úcovýmifaktormi,ktorénegatívnedeterminujúQoL,súperzistujúcabolest’,pato- logickáúnavaarannástuhnutost’.Bolest’,prítomnávpokojiajpriaktivite,narúšaspánkovú architektúruavedieksyndrómuchronickéhovycerpania(fatigue),ktorýpacienticastohod- notia ako viac limitujúci než bolest’ samotnú. Ranná stuhnutost’, trvajúca nezriedka hodiny, blokujerannérutinnécinnosti.Progresívnastrataklbovýchrozsahovasvalovejsilyznemož- nujesebaobsluhu(hygiena,obliekanie),covediekzávislostinapomocidruhejosoby,strate autonómieapocitombezmocnosti.Vznikátakcirculusvitiosus:bolest’--inaktivita--atrofia --zhoršeniefunkcie(Krchnaváetal.,2018).


\subsection{Psychologickádoména}

Chronický, celoživotný charakter choroby s nepredvídatel’ným striedaním relapsov a remisií, spolu s hrozbou trvalého postihnutia, generuje u pacientov chronickú úzkost’ a depresívne poruchy. Prevalencia depresie sa odhaduje na 15 -- 40\%, pricom je známe, že depresia znižuje prah bolesti a zhoršuje adherenciu k liecbe. Významným faktorom je aj alterácia telesného obrazu (body image) v dôsledku viditel’ných deformít, co vedie k so- ciálnemustiahnutiu.Pacienticastozažívajúfenomén"naucenejbezmocnostiästratuidentity (Olejárová;kol.,2016).


\subsection{Sociálnaaekonomickádoména}

RA postihuje jedincov na vrchole produktívneho veku, co vedie k vysokej miere in- validizácie a strate pracovnej schopnosti. Absencia príjmu prináša ekonomickú instabilitu. Znížená mobilita vedie k sociálnej izolácii a zanechaniu zál’ub. Ochorenie mení dynamiku rodinných vzt’ahov, ked’ sú partneri nútení prevziat’ rolu opatrovatel’ov. Úlohou fyziotera- pie je prostredníctvom edukácie a tréningu kompenzacných mechanizmov maximalizovat’ funkcnúnezávislost’atýmzlepšit’QoL(Krchnaváetal.,2018;Poliakováetal.,2019). 12


\section{Diferenciálnadiagnostika}

Proces diferenciálnej diagnostiky pri RA je intelektuálnou výzvou, ktorá vyžaduje rigorózne vylúcenie iných nozologických jednotiek s prekrývajúcou sa symptomatológiou. Správna a vcasná diferenciácia je sine qua non pre vol’bu adekvátnej terapie, nakol’ko lie- cebnéprístupysaprijednotlivýchochoreniachdiametrálnelíšia(KarelPavelkaetal.,2018).


\subsection{Osteoartróza(OA)}

Osteoartrózajenajcastejšímïmitátorom"RA,predovšetkýmugeriatrickejpopulácie. Na rozdiel od RA, ktorá je primárne zápalovým ochorením, je OA degeneratívny proces spojenýsopotrebovanímchrupavky. • Charakter bolesti: Pri OA je bolest’ mechanická -- zhoršuje sa námahou a v pokoji ustupuje.PriRAjebolest’zápalová,prítomnáajvkl’ude. • Stuhnutost’:PriOAjekrátka,"štartovacia"(do30minút).PriRAjedlhodobá. • Lokalizácia:OApostihujetypickyDIPklby(Heberdenoveuzly),PIPklby(Bouchar- doveuzly)aklbpalcaruky(rhizartróza).RApostihujeMCPklbyazápästie. • Laboratórium: Pri OA sú zápalové markery (CRP, FW) v norme (Olejárová; kol., 2016).


\subsection{Séronegatívnespondyloartritídy}

Ideoskupinuzápalovýchochorení,prektoréjetypickáneprítomnost’reumatoidného faktoraaväzbanaantigénHLA-B27. 1. Psoriatickáartritída(PsA):Môžebyt’klinickyvel’mipodobnáRA.Typickýmidife- renciacnýmiznakmisúprítomnost’psoriázynakoži(hociartritídamôžepredchádzat’ kožnýmprejavom),postihnutieDIPklbov,daktylitída(párkovitýopuchceléhoprsta)a zmenynanechtoch(jamky,onycholýza).NaRTGvidímesúcasneerózieajnovotvorbu kosti(Macejová,2019). 2. Ankylózujúca spondylitída (Bechterevova choroba): Dominuje postihnutie axiál- neho skeletu -- zápalová bolest’ chrbta a sakroiliitída. Periférna artritída je zvycajne asymetrickáatýkasavel’kýchklbovdolnýchkoncatín(Gravaldietal.,2022). 3. Reaktívna artritída: Vzniká v casovej súvislosti s prekonanou urogenitálnou alebo crevnouinfekciou.KlasickáReiterovatriádazahrnaartritídu,uretritíduakonjunktivi- tídu(Schereretal.,2020).


\subsection{Kryštálovéartropatie(Dna)}

Dnavá artritída je metabolické ochorenie spôsobené depozíciou kryštálov nátrium- urátuvklboch. 13 • Priebeh: Typický je akútny, perakútny dnavý záchvat s extrémne silnou bolest’ou, opuchom a zacervenaním, ktorý vzniká náhle, casto v noci ("podagra"-- postihnutie MTPklbupalcanohy). • Chronickáforma:Tofóznadnamôžeimitovat’RAprítomnost’ouuzlov(tofov)aeró- zií, ale rozlíšenie umožní punkcia synoviálnej tekutiny a dôkaz kryštálov (Rovenský, 2001).


\subsection{Systémovýlupuserythematosus(SLE)}

SLEjemultisystémovéautoimunitnéochorenie. • Klby: Artritída pri SLE býva polyartikulárna a symetrická ako pri RA, ale je typicky neerozívna(Jaccoudovaartropatia)--nenicíkost’,hocispôsobujedeformity. • Ostatnépríznaky:Fotosenzitivita,motýl’ovitýexantémnatvári,postihnutieobliciek (nefritída), serozitídy. Typická je prítomnost’ anti-dsDNA a anti-Sm protilátok (Ro- venskýetal.,2015).


\section{Terapiareumatoidnejartritídy}

Manažment RA prešiel v posledných dvoch dekádach revolucnými zmenami. Sú- casná stratégia je založená na koncepte "Treat to Target" (liecba k ciel’u), ktorého im- peratívom je dosiahnutie klinickej remisie alebo aspon stavu nízkej aktivity choroby v co najkratšom case. Tento prístup si vyžaduje proaktívnu úpravu liecby a úzku spoluprácu re- umatológa,pacientaarehabilitacnéhotímu(Smolenetal.,2018).


\subsection{Farmakologickáliecba}

Farmakoterapia zostáva pilierom liecby a delí sa na symptomatickú, ktorá len tlmí prejavy,achorobumodifikujúcu,ktorámenípriebehochorenia. Symptomatickáliecba(NSAaKortikoidy) Tieto lieky poskytujú pacientovi rýchlu úl’avu, ale neriešia podstatu -- nezastavujú klbovúdeštrukciu. • Nesteroidné antiflogistiká (NSA): Efektívne tlmia bolest’ a rannú stuhnutost’ inhi- bíciou enzýmov cyklooxygenázy. Ich dlhodobé užívanie je však limitované rizikom gastropatie(vredy),kardiovaskulárnychpríhodarenálnehopoškodenia. • Glukokortikoidy: Majú silný a rýchly protizápalový efekt. Využívajú sa najmä ako premost’ujúca liecba (bridging therapy) vo vcasných štádiách alebo pri vzplanutiach. Pre riziko závažných nežiaducich úcinkov (osteoporóza, diabetes, hypertenzia, kata- rakta) je snahou podávat’ co najnižšiu úcinnú dávku co najkratší cas (Olejárová; kol., 2016). 14 Chorobumodifikujúcelieky(DMARDs) DMARDs (Disease Modifying Antirheumatic Drugs) zasahujú priamo do patogene- tickýchmechanizmovaspomal’ujúrádiografickúprogresiu. 1. KonvencnésyntetickéDMARDs(csDMARDs):Metotrexátje"zlatýmštandardomä liekom prvej vol’by. Pôsobí ako antagonista kyseliny listovej. K d’alším liekom patria leflunomid,sulfasalazínaantimalariká(hydroxychlorochín). 2. Biologické DMARDs (bDMARDs): Sú to geneticky inžinierované proteíny, ktoré cieleneblokujúšpecifickécytokínyalebobunky.Súindikovanéprizlyhaníkonvencnej liecby.Patrísem: • anti-TNFliecba(Adalimumab,Infliximab), • blokátoryIL-6(Tocilizumab), • B-bunkovádeplécia(Rituximab), • modulátorykostimulácieT-buniek(Abatacept). 3. Cielené syntetické DMARDs (tsDMARDs): Ide o tzv. malé molekuly (inhibítory JAK-kináz--tofacitinib,baricitinib),ktorésapodávajúperorálneablokujúvnútrobun- kovýprenossignálu.Súrovnakoúcinnéakobiologiká(Schereretal.,2020).


\section{Komplexnárehabilitácia}

Rehabilitácia je neoddelitel’nou súcast’ou terapeutického plánu. Zatial’ co farmako- terapia kontroluje zápal, rehabilitácia rieši funkcné dopady ochorenia. Jej ciel’om je udržat’ rozsah pohybu, svalovú silu, zmiernit’ bolest’ a naucit’ pacienta žit’ s chorobou tak, aby bol conajdlhšiesebestacný(Kolár;kol.,2012).


\subsection{Kinezioterapiaaliecebnátelesnávýchova}

Pohybová liecba (LTV) musí byt’ prísne individuálna a rešpektovat’ aktuálnu ak- tivitu ochorenia (akútne vs. chronické štádium). Nedávne systematické prehl’ady a meta- analýzy potvrdzujú, že pravidelné cvicenie je bezpecné a má signifikantný pozitívny vplyv na funkcnú kapacitu a aeróbnu zdatnost’ bez zhoršenia aktivity ochorenia alebo poškodenia klbov(Zhangetal.,2025;Metsiosetal.,2020). Cvicenie v akútnom štádiu V case exacerbácie (vzplanutia) zápalu je prioritou ochrana klbu(Kolár;kol.,2012). • Polohovanie: Klby sa polohujú vo funkcne výhodných postaveniach na prevenciu flekcnýchkontraktúr. 15 • Pasívne pohyby: Vykonávajú sa opatrne do prahu bolesti na udržanie klbovej vôl’e a výživychrupavky. • Izometrické cvicenie: Svalové kontrakcie bez zmeny dlžky svalu sú kl’úcové na pre- venciurýchlejsvalovejatrofie,ktorásprevádzazápal. (Kolár;kol.,2012;Navrátiletal.,2019) Cvicenie v chronickom štádiu a remisii V tomto období je ciel’om obnova funkcie a kondície(Metsiosetal.,2020). 1. Aeróbny tréning: Chôdza, cyklistika, plávanie. Odporúca sa stredná intenzita (60 -- 70\% maximálnej tepovej frekvencie) 30 minút denne. Aeróbne cvicenie znižuje kar- diovaskulárne riziko, ktoré je u pacientov s RA zvýšené, a redukuje únavu (Metsios etal.,2020). 2. Silový tréning: Cvicenie proti odporu (s pružnými t’ahmi Thera-Band, l’ahkými cin- kami)jenevyhnutnénabojprotireumatickejkachexii(úbytkusvalovejhmoty).Štúdie ukazujú, že dynamický silový tréning je bezpecný a vedie k zlepšeniu svalovej sily a funkcnejschopnosti(Zhangetal.,2025). 3. Cvicenie ruky: Špecifické posilnovacie a mobilizacné cvicenia pre drobné klby ruky (programSARAH)zlepšujúsiluúchopuamanuálnuzrucnost’efektívnejšieakobežná starostlivost’(Williamsetal.,2015). High-intensity interval training Vysokointenzívny intervalový tréning (HIIT) predsta- vuje inovatívny prístup, ktorý bol dlho považovaný za nevhodný pre pacientov so zápalo- vými artropatiami. Nedávne randomizované kontrolované štúdie však priniesli prekvapivé zistenia. Štúdia ExeHeart preukázala, že HIIT v primárnej fyzioterapeutickej starostlivosti je bezpecný a efektívny u pacientov so zápalovou artritídou vrátane RA. Výsledky nazna- cujú zlepšenie aeróbnej kapacity a kardiovaskulárnych parametrov bez zhoršenia zápalovej aktivityochorenia(Nordénetal.,2024). Kl’úcom k úspešnej implementácii HIIT u pacientov s RA je individuálna adaptá- cia protokolu. Intervaly vysokej intenzity by mali byt’ kratšie a odpocinkové fázy dlhšie v porovnaní so zdravou populáciou. Cvicenie by malo prebiehat’ na nebolestivých klboch s preferenciou aktivít s nízkou nárazovou zát’ažou, ako je cyklistika alebo plávanie (Nordén etal.,2024;Metsiosetal.,2020). Podporazmenyživotnéhoštýlu Fyzioterapeutomvedenéintervenciezameranénazmenu správania (behavior change interventions) sa ukazujú ako sl’ubný prístup na zvýšenie dlho- dobej adherencie k pohybovej aktivite. Štúdia PIPPRA preukázala realizovatel’nost’ takejto 16 intervencieupacientovsRA.Kombináciaedukácie,stanoveniaindividuálnychciel’ovapra- videlnéhofollow-upuvediekudržatel’némuzvýšeniufyzickejaktivity(Larkinetal.,2024).


\subsection{Efektívnost’fyzioterapiepriRA}

Systematické prehl’ady a meta-analýzy poskytujú vysokú úroven dôkazov pre úcin- nost’ fyzioterapie v manažmente RA. Najnovší systematický prehl’ad analyzujúci randomi- zované kontrolované štúdie dospel k záveru, že fyzioterapeutická intervencia vedie k sig- nifikantnému zmierneniu bolesti a zlepšeniu funkcných parametrov u pacientov s RA (El Ammareetal.,2023). Špecificky pre rehabilitáciu pri reumatických ochoreniach existujú štandardizované protokoly zahrnajúce kombináciu kinezioterapie, fyzikálnych procedúr a edukácie. Kom- plexný rehabilitacný program by mal byt’ zostavený multidisciplinárnym tímom a prispôso- bený individuálnym potrebám pacienta, fáze ochorenia a prítomnosti komorbidít (Jarošová etal.,2018).


\subsection{Fyzikálnaterapia}

Fyzikálna terapia predstavuje integrálnu súcast’ rehabilitacného programu pri RA. Jej základným princípom je využitie fyzikálnych podnetov -- tepla, chladu, elektrických im- pulzov, svetla a mechanickej energie -- na dosiahnutie terapeutických úcinkov v ciel’ových tkanivách. Tieto procedúry primárne slúžia ako príprava pred kinezioterapiou, pricom ich analgezujúci a myorelaxacný efekt umožnuje pacientovi efektívnejšie cvicit’ (Kraft, 2018; Krížová,2014). Termoterapia Aplikácia tepla (hypertermoterapia) nachádza uplatnenie predovšetkým v chronic- kom štádiu ochorenia. Teplo spôsobuje vazodilatáciu, zlepšuje perfúziu tkanív a uvol’nuje svalový spazmus. Najcastejšie využívané modality zahrnajú parafínové zábaly, ktoré sú ob- zvlášt’ vhodné pre drobné klby rúk, a solux lamp. V akútnom štádiu zápalu je naopak in- dikovaná kryoterapia (chlad), ktorá tlmí zápalovú reakciu a znižuje opuch (Vavrová et al., 2016). Elektroliecba Elektroterapeutické modality majú v liecbe RA dlhodobú tradíciu. TENS (transku- tánna elektrická nervová stimulácia) poskytuje analgéziu prostredníctvom vrátkového me- chanizmutlmeniabolesti.Interferencnéprúdy(IFP)kombinujúvýhodystredofrekvencných a nízkofrekvencných prúdov a sú suverénnou modalitou pre hlboko uložené klby. Podrob- nejšie o IFP pojednáva samostatná kapitola tejto práce (Podebradský et al., 2009; Navrátil etal.,2019). 17 Fototerapiaamechanoterapia Nízkohladinovýlaser(LLLT)aultrazvukpodporujúhojivéprocesyvmäkkýchtkani- vách prostredníctvom fotobiomodulacných a termických úcinkov. Tieto modality sú vhodné pri tendosynovitíde, entezopatii a na podporu hojenia po artikulárnych injekciách (Podeb- radskýetal.,2009).


\subsection{Balneoterapiaakúpel’náliecba}

Významnou súcast’ou komplexnej starostlivosti o pacientov s RA je balneoterapia, ktorámávnašichpodmienkachdlhorocnútradíciu.Ideovyužitieprírodnýchliecivýchvôd, peloidov (bahna) a klimatických faktorov. Podl’a najnovšej meta-analýzy, ktorú publikovala Aribietal.(2025),prinášabalneoterapiasignifikantnézlepšenievredukciibolestiazlepšení kvality života (QoL) u pacientov s reumatickými ochoreniami. Efekt nie je len placebo; ter- málne minerálne vody pôsobia chemicky (cez kožu sa vstrebávajú minerály ako síra, jód) a fyzikálne (teplo, vztlak). Verhagen et al. (2015) v Cochrane systematickom prehl’ade potvr- dzujú,že"kúpel’náliecba"(Spatherapy)mápozitívnyvplyvnafunkcnúkapacitupacientovs RA,pricombenefitypretrvávajú3až6mesiacovpoukonceníliecebnéhopobytu.Špecifický význam majú sírne vody (napr. Piešt’any, Trencianske Teplice). Štúdia autorov Zapolski et al.(2024)poukázalanato,žesírnekúpelenielentlmiazápal,alemajúajpozitívnyvplyvna kardiovaskulárneparametre,cojeuRApacientov(castotrpiacichaterosklerózou)kl’úcové. Mechanizmus úcinku zahrna pravdepodobne aj moduláciu oxidacného stresu. Wa˛tor (2020) uvádzajú, že balneochemické vlastnosti vôd ovplyvnujú redoxný potenciál v tkanivách, cím napomáhajúneutralizáciivol’nýchradikálovvznikajúcichprichronickomzápale.


\subsection{Ergoterapiaaochranaklbov}

Ergoterapia ucí pacienta techniky "Joint Protection"-- ako vykonávat’ bežné cinnosti bez zbytocného pret’ažovania klbov. Zahrna používanie dlahovania na prevenciu deformít, úpravu domáceho prostredia a používanie kompenzacných pomôcok (napr. špeciálne otvá- race,zosilnené rúckypríborov), ktoréšetria drobnéklbyruky aumožnujú pacientovi zostat’ nezávislým(Herbertetal.,2022). Rehabilitacné ošetrovatel’stvo doplna ergoterapeutickú intervenciu o nácvik základ- nýchsebaobslužnýchaktivítupacientovst’ažkýmiporuchamihybnosti.Ciel’omjemaxima- lizácia funkcnej nezávislosti a prevencia sekundárnych komplikácií z imobility (Klusonová etal.,2014).


\subsection{Rehabilitáciaareumatológia--interdisciplinárnaspolupráca}

Optimálny manažment pacienta s RA vyžaduje koordinovanú spoluprácu reumato- lóga, fyzioterapeuta, ergoterapeuta, rehabilitacného lekára a d’alších špecialistov. Rehabi- litacná intervencia by mala byt’ integrovaná do liecebného plánu od momentu stanovenia diagnózyanieažporozvojiireverzibilnýchdeformít(Gúth,2010). 18 Vcasná rehabilitácia zahrna nielen kinezioterapiu a fyzikálnu liecbu, ale aj edukáciu pacienta o ochorení, nácvik kompenzacných stratégií a psychologickú podporu. Pacient by mal byt’ aktívnym partnerom v liecebnom procese a mat’ realistické ocakávania ohl’adom priebehuochoreniaamožnostíterapie(Gúth,2010;Špiritovic,2018).


\section{Prognózareumatoidnejartritídy}

PrognózaRAjevysokovariabilnáazávisíodmnožstvafaktorov.Identifikáciaprog- nostických markerov umožnuje stratifikovat’ pacientov podl’a rizika a prispôsobit’ intenzitu liecby(Wolfeetal.,2010).


\subsection{Faktoryovplyvnujúcedlhodobývýsledok}

Nepriaznivéprognostickéfaktoryzahrnajúvysokétitrereumatoidnéhofaktoraaanti- CCP protilátok, prítomnost’ erozívnych zmien v skorom štádiu ochorenia, extraartikulárne manifestácie,vysokúzápalovúaktivituafunkcnúdisabilitunazaciatkuchoroby.Ženymajú všeobecne horšiu prognózu než muži a genetické faktory (HLA-DR4) sú spojené s agresív- nejšímpriebehom(Wolfeetal.,2010). Vcasné zahájenie liecby DMARDs, dosiahnutie remisie do 6 mesiacov od zaciatku symptómov a striktná kontrola zápalovej aktivity ("Treat to Target") signifikantne zlepšujú dlhodobé výsledky. Pacienti, ktorí dosiahnu remisiu v ökne príležitosti", majú dobré šance na dlhodobé funkcné zachovanie bez významného štrukturálneho poškodenia (K. Pavelka, 2012;McCarthyetal.,2018).


\subsection{Mortalitaakvalitaživota}

Napriek pokrokom v terapii zostáva stredná dlžka života pacientov s RA skrátená v porovnaní s bežnou populáciou, a to približne o 5 až 10 rokov. Hlavnými prícinami smrti sú kardiovaskulárne ochorenia, infekcie a malignity. Moderná agresívna terapia však viedla k zlepšeniuprežívaniavposlednýchdekádach(Wolfeetal.,2010). Kvalita života súvisí predovšetkým s funkcným stavom a kontrolou bolesti. Pacienti v remisii alebo s nízkou aktivitou ochorenia udávajú kvalitu života porovnatel’nú s bežnou populáciou. Súcasná terapeutická stratégia preto kladie dôraz nielen na kontrolu zápalu, ale ajnamaximalizáciufunkcnejkapacityaaktívnezapojeniepacientadobežnéhoživota(Ole- járová,2008). 19 2 Interferencné prúdy Interferencnáterapia(IFP),castooznacovanáako"král’ovnáelektroliecby",predsta- vuje jednu z najsofistikovanejších modalít vo fyzikálnej medicíne. Patrí do skupiny stredo- frekvencných prúdov, ktoré vd’aka svojim unikátnym fyzikálnym vlastnostiam dokážu pre- konat’ bariéru kože s minimálnym odporom a dorucit’ tak terapeuticky úcinnú elektrickú energiu priamo do hlboko uložených tkanív -- svalov, nervov a klbov. Táto kapitola ponúka hlbkovúanalýzuteoretickýchvýchodísk,biofyzikálnychinterakciíaklinickýchaplikáciíin- terferencných prúdov, pricom osobitný zretel’ je kladený na ich nezastupitel’né miesto v komplexnejrehabilitacnejstarostlivostiopacientovsreumatoidnouartritídou(Capko,1998).


\section{Históriaavývojmetódy}

Genéza tejto revolucnej metódy siaha do polovice 20. storocia a je neodmyslitel’ne spätá s menom rakúskeho fyzika a vynálezcu, Dr. Hansa Nemeca (1907--1981). V 40. ro- koch minulého storocia celila rehabilitácia vážnemu technickému limitu: vtedajšie nízko- frekvencnéprúdy(napr.faradickýimpulzalebodiadynamicképrúdy)malivýbornýliecebný efekt, ale narážali na vysokú impedanciu l’udskej kože. Koža sa pri nízkych frekvenciách správa ako izolátor. Aby sa prúd dostal do hlbky, bolo nutné použit’ vysokú intenzitu, co však vyvolávalo nepríjemné pálenie, bolest’ a v krajných prípadoch až poleptanie pokožky podelektródou.Pacientitútoliecbucastonetolerovali(Capko,1998). Dr. Nemec prišiel s geniálnym riešením, ktoré obišlo tento problém. Všimol si, že s rastúcou frekvenciou klesá odpor kože. Navrhol preto použit’ dva nezávislé prúdové okruhy so strednou frekvenciou (okolo 4000 Hz), pre ktoré je koža takmer priehl’adná. Kl’úcovou myšlienkou bolo, že tieto dva prúdy sa "stretnúäž v hlbke tkaniva. V mieste ich prekríženia dôjde k fyzikálnemu javu interferencie, cím vznikne nový, nízkofrekvencný pulz priamo v ciel’ovom orgáne. Tento princíp ëndogénneho generovania"liecivého signálu znamenal zá- sadnýzlom.Kýmprvéprístrojeznámeako"Nemectron"využívalistatickúinterferenciu,ne- skoršítechnologickývývojumožnilvznikdynamickýchvektorovýchpolí,schopnýchošetrit’ ajrozsiahleaanatomickyzložitéoblasti(Podebradskýetal.,2009;Goats,1990).


\section{Fyzikálneprincípyinterferencie}

Fundamentálnym princípom metódy je superpozícia vlnení. Do organizmu sú pri- vádzané dva sínusové striedavé prúdy strednej frekvencie pomocou štyroch elektród. Prvý okruh má konštantnú nosnú frekvenciu f (napríklad 4000 Hz) a druhý okruh má frekven- 1 ciu f mierneodlišnú(napríklad4100Hz).Tietoprúdyprechádzajútkanivombeztoho,aby 2 pacientadráždili.Vmieste,kdesaichdráhypretnú(geometrickýstredpolohyelektród),do- chádza k ich vzájomnému skladaniu -- interferencii. Výsledkom je vznik amplitúdovej mo- dulácie. Frekvencia tejto modulácie nie je náhodná; presne zodpovedá rozdielu frekvencií 20 oboch nosných prúdov. Tento parameter nazývame Amplitúdovo modulovaná frekvencia (AMF): AMF =|f − f |=|4100−4000|=100Hz 2 1 Zklinickéhohl’adiskatoznamená,žehocitkanivompretekáprúd4000Hz(ktorýnecítime), bunkyvhlbkereagujúna"rytmus"100Hz(ktorýlieci).Týmtospôsobomdokážemedorucit’ úcinok nízkej frekvencie (1--100 Hz) do hlbky, kam by sa priama nízka frekvencia cez kožu nedostala(Robertsonetal.,2006;Goats,1990).


\subsection{Fenoménprekonaniakožnéhoodporu}

Elektrické vlastnosti kože sú kl’úcom k pochopeniu výhody IFP. Koža sa správa ako kapacitnýodpor(kondenzátor).Jejimpedancia(Z)jenepriamoúmernáfrekvenciiprúdu(f): 1 Z ≈ f Pri klasickej frekvencii 50 Hz je odpor kože enormný, približne 3200 Ω. To predstavuje bariéru, ktorá bráni prieniku energie. Avšak pri použití "nemecovskej"frekvencie 4000 Hz tento odpor klesá strmo nadol, až na hodnotu okolo 40 Ω. Prúd tak vstupuje do tela l’ahko, "hladko", bez dráždenia vol’ných nervových zakoncení v koži (nociceptorov). Pacient necíti žiadne pálenie povrchu, ale vníma príjemné, masívne vibrácie v hlbke svalu alebo klbu. To umožnuje terapeutovi použit’ podstatne vyššie intenzity prúdu než pri iných metódach, co vedieksilnejšiemubiologickémuúcinku(Navrátiletal.,2019;Watson,2020).


\section{Fyziologickéúcinkynaorganizmus}

SpektrumúcinkovIFPjedeterminovanépredovšetkýmvol’boufrekvencieAMF.Re- umatológiavyužívanajmätietotrihlavnéefekty:


\subsection{Analgetickýúcinok(Tlmeniebolesti)}

Potlaceniebolestijenajcastejšouindikáciou.Mechanizmusjeduálnyazávisíodfrek- vencie: 1. Vrátková teória bolesti (Gate Control Theory): Aplikácia vyšších frekvencií AMF (80 -- 100 Hz) silne stimuluje hrubé myelinizované nervové vlákna typu Aβ, ktoré vedú dotykové a tlakové informácie. Tieto rýchle vzruchy prichádzajú do zadných rohovmiechyskôrnežsignálybolestivedenépomalýmivláknamiC.Týmsaaktivujú inhibicné interneuróny, ktoré "zatvoria vrátka"pre bolest’. Efekt je okamžitý, pacient cítiúl’avupriamopocasprocedúry(Kolár;kol.,2012). 2. Opioidnýmechanizmus:Prinízkych,rytmickýchfrekvenciách(1--10Hz)dochádza k stimulácii produkcie endogénnych opiátov (endorfínov a enkefalínov) v centrálnom nervovom systéme. Tento typ analgézie nastupuje pomalšie (po cca 15-20 minútach), 21 ale pretrváva dlhodobo, casto niekol’ko hodín po ukoncení terapie (Navrátil et al., 2019).


\subsection{Myorelaxacnýúcinok(Uvol’neniesvalov)}

Pri reumatoidnej artritíde sú svaly v okolí postihnutého klbu casto v reflexnom krci (spazme),cozvyšujebolest’aobmedzujepohyb.IFPvyužívajavzvanýWedenskéhoútlm. Pri nosných frekvenciách okolo 4000 Hz nervové vlákno nie je schopné reagovat’ na každý jeden impulz depolarizáciou, pretože sa dostáva do stavu únavy (refraktérnej fázy). Dochá- dza k zníženiu dráždivosti a následnému uvol’neniu svalového tonusu (Podebradský et al., 2009).


\subsection{Vazodilatacnýaantiedematóznyúcinok}

Zlepšenieprekrveniajevitálnepreodvádzaniezápalovýchmetabolitov(laktát,pros- taglandíny, histamín). IFP tlmí aktivitu sympatického nervového systému, co vedie k rela- xácii hladkej svaloviny v stenách ciev a ich rozšíreniu (vazodilatácii). Súcasne, rytmické striedanie kontrakcie a relaxácie svalov pri nižších frekvenciách (10-50 Hz) funguje ako "svalová pumpa", ktorá mechanicky vytláca venóznu krv a lymfu z oblasti, cím efektívne redukujeopuchy(edémy)(Capko,1998).


\section{Metodikaaplikácieavariácie}

Správnatechnikajepodmienkouúspechu.Vpraxirozlišujemedvazákladnéspôsoby aplikácie,ktorésalíšiapoctomelektródamiestomvznikuinterferencie.


\subsection{Dvojpólováinterferencia(Premodulovaná)}

Tento režim využíva len dve elektródy (jeden kanál). Špecifikom je, že k zmiešaniu prúdov a vzniku nízkofrekvencnej obálky nedochádza v tele pacienta, ale už v samotnom prístroji. Z elektród teda vystupuje už hotový, modulovaný prúd. Výhody: Je to ideálna technika pre ošetrenie malých plôch, trigger bodov alebo drobných klbov ruky, kam by sa štyri elektródy fyzicky nezmestili. Nevýhody: Nakol’ko prúd je modulovaný už pred vstu- pomdokože,ciastocnesastrácavýhodal’ahkéhoprienikutypickápreklasickúinterferenciu. Hlbkovýefektjepretooniecoslabší(Watson,2020).


\subsection{Štvorpólováinterferencia(Klasická)}

Zlatý štandard pri liecbe vel’kých klbov a chrbtice. Vyžaduje použitie štyroch elek- tród zapojených do dvoch krížiacich sa okruhov. Liecivý úcinok vzniká presne v strede, v tzv.zóneinterferencie.Modernéprístrojeumožnujúmodifikovat’tvartohtopol’a: • Statické pole: Zóna úcinku je fixná, má tvar štvorlístka. Vhodné pre lokalizovanú bolest’. 22 • Vektorové pole (Dynamické): Prístroj cyklicky mení intenzitu jedného okruhu, cím dochádza k "rotácii"pol’a. Interferencná zóna sa tak pohybuje a premasíruje celú ob- last’ medzi elektródami. Toto je metóda vol’by pri difúznych bolestiach klbov (napr. celérameno,koleno)(Navrátiletal.,2019).


\subsection{Aplikáciasvákuovoujednotkou(Vákuováterapia)}

Vel’mi casto sa v klinickej praxi stretneme s kombináciou interferencných prúdov a vákuovej jednotky. Elektródy nie sú pripevnené elastickými pásmi, ale sú umiestnené vo vnútri špeciálnych pohárikov (baniek), ktoré sa prisajú na kožu pomocou podtlaku. Táto kombináciaprinášasynergickýefektdvochterapiívjednomcase: 1. Zlepšenie kontaktu: Podtlak zabezpecí dokonalé pril’nutie elektródy na nerovný po- vrchtela(napr.rameno,lopatka),cozarucujekonštantnýprenosprúdu. 2. Hyperémia(prekrvenie):Samotnýpodtlakvtiahnekožuapodkožiedopohárika,cím spôsobímechanickúvazodilatáciu.Zlepšenéprekrveniepodelektródouznižujekožný odporeštevýraznejšienežsamotnáfrekvencia4000Hz,címsazvyšujeúcinnost’pre- nosuenergiedohlbky. 3. Mikromasáž: Moderné vákuové jednotky pracujú v pulznom režime, kedy sa sila prisatia rytmicky mení. Táto mechanická pulsácia pôsobí ako jemná vákuová masáž, ktorástimulujelymfatickýobehapomáhapriredukciiopuchov. Pri použití vákuovej jednotky je však nutné dbat’ na opatrnost’ u pacientov s krehkými cie- vami (varixy) alebo u seniorov s atrofickou kožou, kde by silný podtlak mohol spôsobit’ petéchiealebohematómy(Podebradskýetal.,2009).


\section{IndikácieakontraindikáciepriRA}

\subsection{Indikácie}

IFP je suverénnou metódou vol’by najmä v subakútnom a chronickom štádiu RA. Hlavnýmiciel’misú: 1. Analgesia:Zvládnutiebolestibeznutnostizvyšovat’dávkyliekov. 2. Resorpcia:Urýchlenievstrebávaniachronickýchburzitídasynoviálnychvýpotkov. 3. Relaxácia:Uvol’neniereflexnýchspazmovparavertebrálnehosvalstvaasvalovvokolí postihnutýchklbov. 4. Trofika: Zlepšenie výživy tkanív pri akrocyanóze a Raynaudovom syndróme (Kolár; kol.,2012;Capko,1998). 23


\subsection{Kontraindikácie}

HocijeIFPbezpecná,existujústavyvylucujúcejejpoužitie: • Absolútne: Gravidita (aplikácia na driek a panvu), kardiostimulátor (riziko elektro- magnetickej interferencie), febrilné stavy, akútna tromboflebitída a flebotrombóza (ri- zikoembolizácie!),onkologickéochoreniavmiesteaplikácie. • Relatívne/Lokálne: Aplikácia v oblasti ocí, karotických sínusov. Poruchy citlivosti kože vyžadujú opatrnost’. Prítomnost’ kovových implantátov (TEP) nie je absolútnou kontraindikáciou, nakol’ko stredofrekvencný prúd nemá galvanickú zložku a nespô- sobuje chemické popálenie, avšak odporúca sa nižšia intenzita kvôli riziku prehriatia kovu(Watson,2020).


\section{PorovnanieIFPsinýminízkofrekvencnýmiprúdmi}

Pre správnu indikáciu je kl’úcové chápat’, v com sa interferencné prúdy líšia od kla- sických nízkofrekvencných metód, ako sú TENS (Transkutánna elektrická nervová stimulá- cia)alebodiadynamicképrúdy(DD).


\subsection{Interferenciavs.TENS}

TENS je najcastejšie používanou metódou domácej liecby bolesti. Pracuje s nízkou frekvenciou(1-150Hz)priamonakoži. • TENS: Má silný povrchový úcinok, ale t’ažšie preniká do hlbky klbu kvôli odporu kože.Pacientcítiostré"bodanieälebomravceniepriamopodelektródou. • IFP:Vd’akanosnejfrekvencii4000Hzprekonávakožnýodporl’ahšie.Úcinokvzniká ažvhlbke(ëndogénne").PacientitolerujúvyššieintenzityprúdupriIFPnežpriTENS, covediekhlbšiemumyorelaxacnémuúcinku(Goats,1990).


\subsection{Interferenciavs.Diadynamicképrúdy(DD)}

DDprúdy(podl’aBernarda)sústaršoumetódouvyužívajúcousínusovéprúdy50/100 Hz. • Tolerancia:DDprúdymajúsilnúgalvanickúzložku(jednosmernú),ktorádráždikožu aspôsobujehyperémiu(zacervenanie)podelektródou.Hrozírizikopoleptaniakožepri dlhšejaplikácii. • Komfort: IFP nemajú galvanickú zložku (sú cisto striedavé), preto nespôsobujú che- mické zmeny pod elektródou ani pálenie. Sú bezpecnejšie pre pacientov s citlivou po- kožkou a môžu sa aplikovat’ aj v prítomnosti kovových implantátov (s opatrnost’ou), kýmDDsúprikovochprísnekontraindikované(Podebradskýetal.,2009). 24


\section{Súcasnýstavdôkazov(Evidence-BasedMedicine)}

Véremedicínyzaloženejnadôkazochjenevyhnutnéoverovat’empirickéskúsenosti exaktnými dátami z klinických štúdií. Výskum úcinnosti interferencných prúdov v posled- nýchrokochpriniesolviacerovýznamnýchzistení,ktorépotvrdzujúichmiestovmoderných liecebnýchprotokoloch.


\subsection{Úcinnost’pribolestiachchrbticeaklbov}

Najnovší systematický prehl’ad a meta-analýza, ktorú publikovali Érika P Rampazo et al. (2023), sa zamerala na pacientov s chronickou nešpecifickou bolest’ou dolnej casti chrbta (LBP), co je castá komorbidita u pacientov s reumatickými ochoreniami. Autori ana- lyzovalidátazviacerýchrandomizovanýchkontrolovanýchštúdiíadospelikzáveru,žeIFP má štatisticky významný efekt na zníženie intenzity bolesti oproti placebu. Dôležitým ziste- ním bolo, že IFP nielenže znižuje bolest’ bezprostredne po aplikácii, ale zlepšuje aj funkcné parametre(pohyblivost’),cojeprerehabilitáciukl’úcové. Podobnepovzbudivévýsledkyprinieslaštúdiazameranánaosteoartrózukolena,ktorú viedliChenetal.(2022).Hocisazameralinagonartrózu,patofyziologickýmechanizmusbo- lestijepodobnýakopriRA.Ichmeta-analýzapotvrdila,žeinterferencnáterapiajeefektívna nielen v manažmente bolesti, ale významne prispieva k zlepšeniu kvality života pacientov. V porovnaní s inými elektroterapeutickými modalitami sa IFP javí ako superiorná najmä vd’aka schopnosti preniknút’ do hlbky klbneho puzdra, co povrchové prúdy (napr. TENS) nedokážuvtakejmiere.


\subsection{Neuromodulacnýpotenciál}

Zaujímavý pohl’ad prináša práca Moore et al. (2018), ktorá skúmala využitie IFP pri neuromodulácii gastrointestinálnych porúch. Hoci to priamo nesúvisí s liecbou klbov, potvrdzuje to schopnost’ interferencných prúdov ovplyvnovat’ autonómny nervový systém. Tentopoznatokjeaplikovatel’nýajvreumatológii,kdemoduláciasympatika(vazodilatacný efekt) zohráva úlohu pri liecbe trofických zmien a zlepšení mikrocirkulácie v zapálenom tkanive. Všeobecnemožnokonštatovat’,žeIFPpredstavujebezpecnú,neinvazívnualternatívu k farmakologickej liecbe bolesti, cím umožnuje znížit’ spotrebu analgetík a nesteroidných antireumatík,ktorémajúpridlhodobomužívanívážnevedl’ajšieúcinky(Goats,1990).


\section{PraktickáaplikáciapriRA}

Aplikácia interferencných prúdov u pacientov s reumatoidnou artritídou vyžaduje zohl’adnenie špecifík tohto ochorenia. Na rozdiel od degeneratívnych artropatií, kde je zá- palová zložka sekundárna, pri RA dominuje primárny imunitne mediovaný zápal. Preto je vol’ba parametrov stimulácie a nacasovanie terapie kl’úcová pre dosiahnutie optimálneho terapeutickéhoefektu. 25


\subsection{Parametrestimulácie}

1. Akútna bolest’ a edém: Volíme vyššie frekvencie AMF, typicky 90 -- 100 Hz, prí- padnesmalýmrozkmitom(Spectrum)napr.10Hz.Ciel’omjeanalgéziacezvrátkovú teóriu. Intenzita musí byt’ striktne podprahová (jemné mravcenie), nesmie vyvolat’ motorickúreakciu.Dobaaplikáciekratšia,10--15minút. 2. Chronickástuhnutost’svalov:VolímenízkefrekvencieAMF,1--10Hz,alebostrie- davérežimy(1--100Hz).Ciel’omjemyorelaxáciaa"gymnastikaciev".Intenzitamôže byt’vyššia,ažnahranicumotorickéhoprahu(viditel’néchveniesvalu).Dobaaplikácie dlhšia,15--20minút. 3. Zlepšenietrofikytkanív:PrivaskulárnychkomplikáciáchRA(Raynaudovfenomén, akrocyanóza)volímestrednúfrekvenciu50Hzsvektorovýmpol’om.Kombinujemes podtlakovýmielektródamipresynergickývazodilatacnýefekt(Capko,1998).


\subsection{Aplikacnéschémyprejednotlivéklby}

Ramenný klb: Využíva sa tetrapolárna aplikácia s vektorovým pol’om. Jedna elek- tróda sa umiestnuje ventrálne, druhá na úpon deltoidu, tretia a štvrtá dorzálne. Interferencné pole tak prestúpi celým klbnym puzdrom. Ideálne je použitie vákuových elektród pre lepší kontakt na oblých tvaroch ramena. Pri frozen shoulder (adhezívna kapsulitída), ktorá môže komplikovat’RA,volímekombinovanýprotokol:najprvanalgézia(100Hz,10min),potom myorelaxácia(1-10Hz,10min)(Navrátiletal.,2019). Kolenný klb: Klasická aplikácia "do kríža". Elektródy obklopujú patelu. Pri osla- bení kvadricepsu volíme frekvencie okolo 50 Hz, ktoré majú tonizacný vplyv na svalstvo. Pri opuchoch volíme 100 Hz na podporu rezorpcie. U pacientov s RA je castá prítomnost’ Bakerovej cysty v popliteálnej oblasti; v takomto prípade je vhodné zahrnút’ túto oblast’ do terapeutickéhopol’a(Kolár;kol.,2012). Ruky a zápästia: Pre ošetrenie drobných klbov rúk, ktoré sú pri RA najcastejšie postihnuté, je suverénne najlepšou metódou "vodný kúpel’". Do dvoch plastových nádob s vodou sa vložia uhlíkové elektródy. Pacient ponorí ruky (alebo nohy) do vody. Voda zabez- pecí dokonalý kontakt s celým nerovným povrchom prstov. Táto metóda umožnuje ošetrit’ všetkyklbyrukynaraz,cojeprilepeníelektródnemožné(Navrátiletal.,2019). Bedrový klb: Hlboká lokalizácia bedrového klbu vyžaduje použitie väcších elek- tród a vyššej intenzity. Tetrapolárna aplikácia s elektródami umiestnenými v oblasti triesla, laterálne od vel’kého trochanteru, gluteálne a posteromediálne. Vektorové pole zabezpecí ošetrenieceléhoklbuvrátanejehohlbokýchštruktúr(Robertsonetal.,2006). Krcná chrbtica: U pacientov s RA je nutná opatrnost’ kvôli možnej atlantoaxiálnej instabilite. Pred aplikáciou IFP na krcnú chrbticu by mal byt’ vylúcený subluxácia C1-C2 pomocou RTG vyšetrenia. Aplikácia sa realizuje paravertebrálne s nízkymi intenzitami a frekvenciami80-100Hzpreanalgéziu(Podebradskýetal.,2009). 26


\subsection{Frekvenciaadlžkaliecby}

Štandardný liecebný cyklus pozostáva z 10-15 sedení aplikovaných denne alebo ob- den.Úcinokjednejprocedúrytrvániekol’kohodín,prisériovejaplikáciidochádzakukumu- lácii efektu. Po ukoncení cyklu nasleduje prestávka 4-6 týždnov pred prípadným opakova- ním.UpacientovschronickouRAmožnoliecbuopakovat’podl’apotreby,typicky2-3cykly rocne(Capko,1998).


\subsection{Kombináciasinýmiterapeutickýmimodalitami}

Interferencné prúdy je možné efektívne kombinovat’ s d’alšími fyzioterapeutickými metódami. Najcastejšie sa IFP aplikuje pred kinezioterapiou -- zníženie bolesti a uvol’nenie svalov umožnuje pacientovi efektívnejšie cvicit’. Kombinácia s hydroterapiou (teplý bazén) je obzvlášt’ výhodná -- vztlak vody odl’ahcuje klby a teplo potencuje vazodilatacný efekt elektrickéhoprúdu(Gúth,2010). Pri lokálnej aplikácii kryoterapie pred IFP u pacientov s akútnym vzplanutím vy- užívameprotizápalovýefektchladunasledovanýanalgetickýmúcinkomelektrickéhoprúdu. Táto kombinácia je preferovaná pred termoterapiou v akútnej fáze zápalu (Podebradský et al.,2009).


\section{Bezpecnostnéaspektyanežiaduceúcinky}

Interferencná terapia patrí medzi najbezpecnejšie elektroterapeutické modality. Stre- dofrekvencný charakter prúdu eliminuje riziko chemického popálenia pod elektródou, ktoré hrozí pri galvanických prúdoch. Napriek tomu je nevyhnutné dodržiavat’ bezpecnostné pro- tokoly(Watson,2020).


\subsection{Možnénežiaduceúcinky}

• Prechodná erytéma: Zacervenanie kože pod elektródami je fyziologickou reakciou navazodilatáciuanevyžadujeprerušenieliecby. • Svalová únava: Pri nadmernej intenzite alebo príliš dlhej aplikácii s motorickou od- poved’oumôžedôjst’kúnavesvalstva. • Nepríjemné pocity: Pacienti môžu pocit’ovat’ mravcenie, vibrácie alebo tlak. Tieto pocitybynemalibyt’bolestivé--vtakomprípadejenutnéznížit’intenzitu. • Vertigo: Pri aplikácii v oblasti krku sa vzácne môže objavit’ závrat’, co vyžaduje okamžitéukoncenieprocedúry(Goats,1990).


\subsection{Monitorovaniepacienta}

Pocasprocedúrybymalterapeutpravidelnekontrolovat’subjektívnepocitypacienta. Otázky by mali smerovat’ k intenzite vnímania prúdu, lokalizácii pocitov a komfortu. U 27 pacientovsreumatoidnouartritídoujepotrebnévenovat’pozornost’ajstavukože--atrofická, ztencenákožapridlhodobejkortikoterapiijenáchylnejšianapoškodenie(Watson,2020).


\section{Záverkapitoly}

Interferencné prúdy predstavujú efektívnu a bezpecnú terapeutickú modalitu s pev- ným miestom v komplexnej rehabilitácii pacientov s reumatoidnou artritídou. Ich unikátna schopnost’ dorucit’ terapeuticky úcinnú nízkofrekvencnú stimuláciu do hlboko uložených tkanív bez dráždenia povrchových štruktúr ich predurcuje pre liecbu klbov postihnutých zá- palovýmprocesom. Súcasná vedecká evidencia podporuje využitie IFP na zvládnutie bolesti, redukciu opuchov a prípravu pacienta na kinezioterapiu. V kombinácii s farmakologickou liecbou a komplexnýmrehabilitacnýmprogramomprispievajúkzlepšeniufunkcnejkapacityakvality životapacientovsRA(ÉrikaPatríciaRampazoetal.,2022;Fuentesetal.,2010). Budúci výskum by mal priniest’ štandardizované protokoly špecificky pre reumato- idnúartritíduadlhodobésledovanieefektivityvkontextemodernejbiologickejliecby. 28 3 Ciel’ práce a metodológia Praktická cast’ diplomovej práce je zameraná na overenie úcinnosti navrhnutej tera- pie.Vtejtokapitoledefinujemehlavnýaciastkovécielepráce,stanovujemehypotézyacha- rakterizujeme výskumný súbor. Dalej popisujeme metodiku zberu dát a postup spracovania získanýchvýsledkov.


\section{Ciel’práce}

Ciel’om diplomovej práce bolo na základe dostupnej literatúry a vlastného výskumu posúdit’úcinnost’fyzioterapiepriliecbepacientovsreumatoidnouartritídou. Hlavný ciel’: Zistit’, ako kombinovaná liecba interferencnými prúdmi a kineziote- rapiou ovplyvnuje intenzitu bolesti, rozsah pohybu a subjektívne hodnotenie kvality života pacientovsreumatoidnouartritídou. Ciastkovéciele: 1. Porovnat’úcinnost’kombinovanejterapie(IFP+kinezioterapia)asamotnejkineziote- rapienaintenzitubolesti. 2. Analyzovat’vplyvkombinovanejterapienarozsahpohybuvpostihnutýchklboch. 3. Zhodnotit’subjektívnevnímaniekvalityživotapacientovpredapoterapii. 4. Porovnat’výsledkyexperimentálnejakontrolnejskupiny.


\section{Hypotézy}

Nazákladestanovenýchciel’ovsmeformulovalinasledujúcehypotézy: H1: Predpokladáme, že pacienti s reumatoidnou artritídou, ktorí absolvujú terapiu interferencnými prúdmi v kombinácii s kinezioterapiou, uvádzajú nižšiu intenzitu bolesti nežpacientiliecenílenkinezioterapiou. H2:Predpokladáme,žekombináciainterferencnýchprúdovakinezioterapiemáväcší pozitívnyvplyvnarozsahpohybuvpostihnutýchklbochakosamotnákinezioterapia. H3:Predpokladáme,žepacientilieceníkombinovanouterapiousubjektívnehodnotia kvalitusvojhoživotapozitívnejšieakopacientivskupinebezelektroterapie.


\section{Charakteristikavýskumnéhosúboru}

Výskumný súbor tvorilo 100 pacientov s potvrdenou diagnózou reumatoidnej artri- tídy. Podmienkou zaradenia do štúdie bola lekárom potvrdená diagnóza RA. Pacientov sme rozdelili do dvoch skupín po 50 jedincov. Experimentálnu skupinu tvorilo 50 pacientov ab- solvujúcich kombinovanú liecbu interferencným prúdom a kinezioterapiou. Kontrolnú sku- pinutvorilo50pacientovpodstupujúcichibakinezioterapiu. 29 Grafzobrazujerozdeleniemužovažienvskupineskombinovanouliecbouajvsku- pinesizolovanoukinezioterapiou. 35,00 30,00 28,00 26,00 24,00 25,00 22,00 20,00 15,00 10,00 5,00 0,00 ITF+kinezioterapia Kinezioterapia votneicaptecoP Rozdeleniesúborupodl’apohlavia Muži Ženy Graf1Pocetmužovažienprejednotlivéskupiny Skupinu s kombinovanou liecbou tvorilo 24 mužov a 26 žien, pricom kontrolná sku- pinazahrnalazastúpenie22mužova28žien.Zhl’adiskarozloženiapohlaviamožnokonšta- tovat’,žeobeskupinybolihomogénne. Okrem základného rozdelenia podl’a pohlavia sme sa zamerali aj na analýzu pra- covného zat’aženia pacientov, ked’že typ zamestnania môže významne ovplyvnovat’ mieru zát’aže pohybového aparátu. V nasledujúcom grafe (Graf 2) je znázornená štruktúra súboru zhl’adiskatypuvykonávanejpráce. 30 35,00 30,00 25,00 21,00 19,00 20,00 18,00 17,00 15,00 13,00 12,00 10,00 5,00 0,00 ITF+kinezioterapia Kinezioterapia votneicaptecoP Rozdeleniesúborupodl’atypuzamestnania Fyzické Sedavé Neaktívne Graf2Rozdeleniepacientovpodl’atypuzamestnania V oboch skupinách bolo zastúpenie typu zamestnania porovnatel’né. V skupine ITF bolo zastúpených 19 pacientov s fyzickým, 18 so sedavým a 13 s neaktívnym životným štýlom. V kontrolnej skupine (kinezioterapia) bolo 17 pacientov s fyzickým, 21 so sedavým a 12 s neaktívnym typom zamestnania. Rozdiely v pracovnom zat’ažení medzi skupinami nebolištatistickyvýznamné. Súcast’ou charakteristiky výskumného súboru bolo aj sledovanie základných demo- grafických a antropometrických parametrov, ako sú vek, telesná výška a hmotnost’. Tieto údaje nám umožnili posúdit’ homogenitu oboch skupín. Ako vyplýva z údajov v Tabul’ke 1 aTabul’ke2,obeskupinyvykazovaliporovnatel’népriemernéhodnotysledovanýchpremen- ných,cosvedcíovysokejmierevyrovnanostivýskumnejvzorky.Tátopodobnost’jedôležitá preobjektivituporovnaniaterapeutickýchvýsledkov. Tabul’ka1DemografickéúdajepacientovvskupineITF+kinezioterapia Pacienti n x¯ sd x min. max. m Vek 50 52,76 5,05 54,5 41 60 Výška 50 173,24 10,13 175,5 159 193 Hmotnost’ 50 73,04 12,35 77,5 60 101 Legenda:n--pocetpacientov,x¯--aritmetickýpriemer,sd--smerodajnáodchýlka,x --medián,min.-- m minimálnahodnota,max.--maximálnahodnota SkupinaITF+kinezioterapiavykazovalapriemernývek52,76±5,05roka,pricompriemernátelesnávýška dosahovala173,24±10,13cmahmotnost’73,04±12,35kg. 31 Tabul’ka2DemografickéúdajepacientovvskupineKinezioterapia Pacienti n x¯ sd x min. max. m Vek 50 51,62 5,90 52,5 40 60 Výška 50 173,04 10,23 174 157 190 Hmotnost’ 50 74,22 11,06 74,0 57 96 Legenda:n--pocetpacientov,x¯--aritmetickýpriemer,sd--smerodajnáodchýlka,x --medián,min.-- m minimálnahodnota,max.--maximálnahodnota Vskupineliecenejsamostatnoukinezioterapioubolpriemernývekpacientov51,62±5,90roka.Priemerná výškatýchtopacientovbola173,04±10,23cmaichpriemernáhmotnost’predstavovala74,22±11,06kg. 32


\section{Metodikavýskumu}

Pre potreby nášho prieskumu sme zvolili metodologický postup kombinujúci sub- jektívnedotazníkovéhodnoteniesobjektívnymmeranímfyzickýchparametrov.Celýproces zberu údajov sa realizoval v prostredí neštátneho zdravotníckeho zariadenia, pricom casový rámec bol stanovený od augusta 2025 do februára 2026. Probandi zapojení do štúdie absol- vovalikompletnýrehabilitacnýprogram.Predjehosamotnýmspustenímsmekládlidôrazna náležité poucenie každého pacienta, ktoré zahrnalo inštruktáž k vyplnaniu formulárov, vy- svetlenie výskumného zámeru a ubezpecenie o prísnej anonymite všetkých spracovávaných údajov. Vstupné vyšetrenie spolu s distribúciou dotazníkov prebehlo v prvý den terapie, za- tial’cokontrolnévýstupnévyšetreniesauskutocnilovzáverecnýden cyklu. Hlavným nástrojom pre monitoring subjektívneho stavu bol nami vytvorený dotaz- ník,vktorompacientikvantifikovalikl’úcovéoblastina10-bodovejstupnici,kdemaximálna hodnotareflektovalanajvyššiumierut’ažkostí.Sledovanéukazovatelesmezameralinainten- zitu klbovej bolesti v pokoji aj pri pohybovej zát’aži, funkcnú mieru sebaobslužných aktivít (ADL), nácvik jemnej motoriky a mobilitu pri chôdzi. Opomenuté nezostali ani typické pri- družené symptómy reumatoidnej artritídy, konkrétne subjektívny pocit únavy a prítomnost’ rannej stuhnutosti. Psychosociálny dopad ochorenia sme mapovali cez kvalitu spánkového cyklu, úroven úzkostných prejavov a vplyv na pracovný výkon. Súbežne so subjektívnym hodnotením sme vykonávali goniometrické merania s ciel’om exaktne stanovit’ aktívny roz- sah pohybu v segmentoch hornej aj dolnej koncatiny. Naša pozornost’ bola upriamená na klby zápästné, drobné klby ruky (MCP a PIP) a klby nohy (MTP), ktoré toto systémové ochoreniepostihujeprimárne. Samotný terapeutický režim bol navrhnutý ako mesacný cyklus pozostávajúci z 12 rehabilitacných jednotiek s frekvenciou trikrát v týždni. Výskumný súbor sme rozdelili na dve porovnávané skupiny s odlišným liecebným protokolom. Experimentálna cast’ súboru podstúpila kombinovanú intervenciu integrujúcu fyzikálnu zložku vo forme interferencných prúdov a cielenú kinezioterapiu. Kontrolná cast’ pacientov absolvovala výhradne pohybovú liecbuvidentickomrozsahu,conámumožnilovalidneposúdit’terapeutickýbenefitelektro- terapieupacientovtrpiacichchronickýmizápalovýmizmenamiklbovéhoaparátu. Aplikáciu interferencných prúdov sme realizovali prostredníctvom prístroja BTL- 4000 Smart. Vzhl’adom na ciel’ dosiahnut’ analgetický efekt a myorelaxáciu sme zvolili tetrapolárnu techniku s konštantne nastavenou frekvenciou 100 Hz, ktorá je v reumatológii preferovaná pre jej výrazný vplyv na tlmenie bolesti. Intenzitu sme u každého probanda la- dili individuálne podl’a subjektívneho vnímania príjemného mravencenia. Každá aplikácia v trvaní 15 minút prebiehala transregionálne v mieste najväcších t’ažkostí. Nadväzujúca ki- nezioterapia v dlžke 20 až 25 minút bola obsahovo identická pre obe skupiny. Intervencia zahrnala techniky mäkkých tkanív, pasívne uvol’novanie a postizometrickú relaxáciu (PIR) svalového aparátu. Aktívna zložka cvicenia bola orientovaná na udržanie klbovej pohybli- vosti, cvicenie úchopových funkcií ruky za pomoci rôznych terapeutických pomôcok a vše- 33 obecnéuvol’neniesdôrazomnaochranuklbovýchplôchpredpret’ažením.


\section{Štatistickéspracovaniedát}

Všetky údajezískané prostredníctvomdotazníka vlastnejkonštrukcie a goniometric- kých meraní sme následne spracovali a vyhodnotili pomocou štatistických metód. Na zá- kladnú charakteristiku súboru a popis sledovaných klinických parametrov (ako sú bolest’, rozsah pohybu a kvalita života) sme aplikovali metódy deskriptívnej štatistiky. Výpocty za- hrnali stanovenie absolútnej pocetnosti, výpocet aritmetického priemeru, smerodajnej od- chýlky, ako aj urcenie minimálnych a maximálnych nameraných hodnôt. Pred samotným testovanímhypotézsmepreverilicharakterrozdeleniadátpomocouShapiro-Wilkovhotestu normality.Vzhl’adomnavýsledky,ktorévoväcšineparametrovnepotvrdilinormálnerozde- leniedát,smepred’alšiuanalýzuzvolilineparametrickétesty. Naporovnanievýsledkovmedziexperimentálnouakontrolnouskupinou(inter-skupinové porovnanie) sme využili neparametrický Mann-Whitneyho U-test. Zmeny v rámci jednot- livých skupín, teda porovnanie hodnôt pred zahájením terapie a po jej ukoncení (intra- skupinové porovnanie), sme posudzovali na základe výsledkov Wilcoxonovho párového testu. Štatistické spracovanie bolo realizované v programe Microsoft Excel. Výstupy sme spracovali do formy prehl’adných tabuliek a grafov (stlpcové a krabicové grafy), pricom hladinaštatistickejvýznamnostibolaprivšetkýchtestochstanovenánahodnotuα =0,05.


\section{Analýzavýsledkov}

Vtejtokapitolepredkladámekompletnéspracovanievýsledkovnášhovýskumu,ktoré vychádza z podrobnej analýzy primárnych dát (100 probandov). Údaje získané z goniomet- rických meraní a subjektívneho dotazníka sme podrobili detailnému spracovaniu. Pre každý sledovaný parameter uvádzame štatistické ukazovatele: priemer (x¯), smerodajnú odchýlku (sd), minimálne (min) a maximálne (max) namerané hodnoty, doplnené o naratívny komen- tárpodl’ametodickýchpožiadaviek.


\subsection{Analýzasubjektívnychparametrov--Bolest’astuhnutost’}

Subjektívne vnímanie klinického stavu je pre pacienta s RA kl’úcové. V tejto casti analyzujemeintenzitubolesti,jejcharakteraproblematikuklbovejstuhnutosti. V experimentálnom súbore (n=50) bola pred terapiou priemerná intenzita bolesti pri zát’aži 7,24 ± 2,10 bodu, pricom minimálna hodnota bola 3 a maximálna 10 bodov (VAS).Poabsolvovaníkomplexnejterapie(ITF+kinezioterapia)kleslapriemernábolest’na 4,92±1,72bodu(min.1,max.9bodov).Bolest’vpokojikleslavtejtoskupinez5,28±2,08 (min.2,max.9)na2,12±1,12bodu(min.1,max.4body). V kontrolnom súbore (n=50) klesla bolest’ pri zát’aži z 7,00 ± 2,06 bodu (min. 3, max.10)na5,40±1,76bodu(min.2,max.10).Bolest’vpokojivkontrolnomsúborepred terapioudosahovalavpriemere5,12±2,05bodu(min.2,max.9)apomesiaciliecbyklesla 34 na3,92±1,80bodu(min.2,max.7bodov).VýsledkyporovnávaTabul’ka3. Tabul’ka3Porovnanieintenzitybolestivpokojiaprizát’aži(body0-10) Parameter/Skupina x¯Pred sdPred min-max x¯Po sdPo min-max Bolest’vpokoji Exp.(ITF+K) 5,28 2,08 2-9 2,12 1,12 1-4 Kontr.(K) 5,12 2,05 2-9 3,92 1,80 2-7 Bolest’prizát’aži Exp.(ITF+K) 7,24 2,10 3-10 4,92 1,72 1-9 Kontr.(K) 7,00 2,06 3-10 5,40 1,76 2-10 10,00 8,00 7,24 7,00 6,00 5,40 4,92 4,00 2,00 0,00 ITF+kinezioterapia Kinezioterapia )01--0(eróksSAV Porovnanieintenzitybolesti(VAS) Predterapiou Poterapii Graf3Porovnanieintenzitybolesti(VAS)predapoterapii Prianalýzecharakterubolestismevexperimentálnejskupinezaznamenalivyrov- nanézastúpenietupej(26%)apulzujúcejbolesti(26%).Nasledovalabolest’bodavá(22%), tlaková (20\%) a pálivá (6\%). V kontrolnej skupine bol pomer síce rôznorodejší, ale do- minovala taktiež tupá bolest’ (28\%), pulzujúca (24\%) a tlaková (22\%). Bodavú bolest’ pocit’ovalo 20\% a pálivú 6\% pacientov. Tieto rozdiely v charaktere bolesti poukazujú na heterogenituvnímaniareumatickýcht’ažkôd. 35 60 50 40 30 28 26 26 24 22 22 20 20 20 10 6 6 0 ITF+kinezioterapia Kinezioterapia )%(einepútsaZ Charakterbolestiuobochskupín Tupá Pálivá Bodavá Pulzujúca Tlaková Graf4Charakterbolestiujednotlivýchskupín S charakterom bolesti úzko súvisí aj casový výskyt bolesti pocas dna. Z analýzy vstupných hodnôt vyplýva, že u oboch skupín dochádza k najväcšej intenzite problémov vrannýchhodinách(ITFn=36,Kontrolnán=34).Popoludnípocit’ovalobolest’14pacientov vITFa16pacientovvkontrolnejskupine. 60 50 40 36 34 30 20 16 14 10 0 Ráno Popoludní )%(einepútsaZ Casovývýskytbolestipocasdna ITF+kinezioterapia Kinezioterapia Graf5Casovývýskytbolestipocasdnauobochskupín 36 Kl’úcovým príznakom je ranná stuhnutost’. V experimentálnom súbore (n=50) merali pacienti intenzitu ranného stuhnutia pred terapiou na úrovni 6,02 ± 2,24 bodu (min. 2, max. 10). Po terapii sa táto hodnota znížila na 3,02 ± 1,95 bodu (min. 1, max. 7). V kon- trolnom súbore (n=50) bol vstupný stav 5,32 ± 2,15 bodu (min. 2, max. 10) and výstupný stav4,42±2,00bodu(min.2,max.9). 10,00 8,00 6,00 6,02 5,32 4,42 4,00 3,02 2,00 0,00 ITF+kinezioterapia Kinezioterapia )01--0(aitunhutsatiznetnI Intenzitarannejstuhnutosti(VAS) Predterapiou Poterapii Graf6Porovnanieintenzityrannejstuhnutosti(VAS) Co sa týka výskytu rannej stuhnutosti, v case vstupného vyšetrenia uvádzalo prí- tomnost’ ranného fenoménu 72\% probandov v experimentálnej skupine (n=36) a 76\% v kontrolnej skupine (n=38). Zvyšná cast’ súboru (28\% resp. 24\%) tento príznak nepo- cit’ovala. Po liecbe v experimentálnej skupine 45\% pacientov uviedlo, že stav stuhnutosti vymizolúplnealebosaobjavujelensporadickypovel’kejzát’aži. 37 100,00 80,00 76,00 72,00 60,00 40,00 28,00 24,00 20,00 0,00 ITF+kinezioterapia Kinezioterapia )%(einepútsaZ Výskytrannejstuhnutostiurespondentov Áno Nie Graf7Výskytrannejstuhnutostiuobochskupín Tabul’ka4PorovnanierannejstuhnutostiaADL(body0-10) Parameter/Skupina x¯Pred sdPred min-max x¯Po sdPo min-max Rannástuhnutost’ Exp.(ITF+K) 6,02 2,24 2-10 3,02 1,95 1-7 Kontr.(K) 5,32 2,15 2-10 4,42 2,00 2-9 ADL Exp.(ITF+K) 5,34 1,85 2-9 3,48 1,62 1-6 Kontr.(K) 5,12 1,28 3-8 4,54 1,66 2-8 Zlepšenievoblastibežnýchdennýchaktivít(ADL)bolovexperimentálnomsúbore markantné(skóreklesloo1,86bodu),cosvedcíoobnovenífunkcnejnezávislostipacientov, zatial’covkontrolnomsúborebolototozlepšenielenopribližne0,5bodu.


\subsection{Analýzagoniometrickýchmeranízápästnéhoklbu}

Prizápästnomklbesmesledovaliflexiuaextenziu.Vobochskupináchdošlokzlep- šeniurozsahupohybu,pricomexperimentálnaskupinavykazovalavýraznejšieprírastky. V experimentálnej skupine (n=50) sme pri flexii v zápästí namerali pred terapiou priemernú hodnotu 60,14 ± 6,94 stupna. Po liecbe sa rozsah zvýšil na 65,50 ± 6,96 stupna. Ukontrolnejskupiny(n=50)bolprogrespriflexiivzápästíz59,54±7,05na61,49±7,14 stupna.TietovýsledkysúznázornenévGrafe8. 38 100,00 80,00 65,50 60,00 60,14 59,54 61,49 40,00 20,00 0,00 Predterapiou Poterapii )enputs(ubyhophaszoR Goniometriazápästia--Flexia ITF+kinezioterapia Kinezioterapia Graf8Goniometriazápästia--Flexia(stupne)uobochskupín Pri extenzii v zápästí sme v experimentálnej skupine (n=50) zaznamenali vstupný priemer50,68±7,46stupnaavýstupnýpriemer55,79±7,47stupna.Vkontrolnejskupine (n=50)došlokzmenez52,18±7,13na54,11±7,16stupna.VýsledkyporovnávaGraf9. 100,00 80,00 60,00 55,79 54,11 50,68 52,18 40,00 20,00 0,00 Predterapiou Poterapii )enputs(ubyhophaszoR Goniometriazápästia--Extenzia ITF+kinezioterapia Kinezioterapia Graf9Goniometriazápästia--Extenzia(stupne)uobochskupín 39 Tabul’ka5Goniometriazápästnéhoklbu--flexiaaextenzia(stupne) Parameter/Strana x¯Pred sdPred min-max x¯Po sdPo min-max Zápästie-Flexia(ExperimentálnaskupinaITF) Dx(vpravo) 60,08 7,29 47-73 65,34 7,45 53-78 Sin(vl’avo) 60,20 6,58 46-72 65,66 6,46 52-77 Zápästie-Flexia(KontrolnáskupinaK) Dx(vpravo) 59,26 7,16 47-72 61,25 7,39 48-74 Sin(vl’avo) 59,82 6,94 50-73 61,73 6,89 50-73 Zápästie-Extenzia(ExperimentálnaskupinaITF) Dx(vpravo) 50,64 7,73 39-67 55,76 7,70 43-70 Sin(vl’avo) 50,72 7,18 38-67 55,82 7,23 42-69 Zápästie-Extenzia(KontrolnáskupinaK) Dx(vpravo) 51,96 7,48 39-68 53,87 7,83 39-67 Sin(vl’avo) 52,40 6,77 38-68 54,35 6,48 40-67


\subsection{Analýzagoniometrickýchmeraníklbovruky(MCPaPIP)}

Pri drobných klboch hornej koncatiny (MCP a PIP) sme u experimentálnej skupiny pozorovalinajvýraznejšiefunkcnéprírastky. V experimentálnom súbore (n=50) sa flexia v MCP klboch zvýšila z vstupných 67,30 ± 7,04 stupna na 72,60 ± 7,30 stupna, co je kl’úcové pre úchopové funkcie. V kon- trolnom súbore sme zaznamenali zmenu z 67,90 ± 8,57 na 69,78 ± 8,20 stupna. Výsledky zachytávaGraf10. 100,00 80,00 72,60 67,30 67,90 69,78 60,00 40,00 20,00 0,00 Predterapiou Poterapii )enputs(ubyhophaszoR GoniometriaMCPklbov--Flexia ITF+kinezioterapia Kinezioterapia Graf10GoniometriaMCP--Flexia(stupne)uobochskupín Extenzia v MCP klboch v experimentálnej skupine (n=50) pred terapiou dosaho- vala34,92±6,55stupnaapomesiaciliecbyvzrástlana39,21±6,31stupna.Vkontrolnej 40 skupine(n=50)smezaznamenaliprogresz34,23±6,82na35,46±6,93stupna(Graf11). 50,00 40,00 39,21 34,92 34,23 35,46 30,00 20,00 10,00 0,00 Predterapiou Poterapii )enputs(ubyhophaszoR GoniometriaMCPklbov--Extenzia ITF+kinezioterapia Kinezioterapia Graf11GoniometriaMCP--Extenzia(stupne)uobochskupín PriflexiivPIPklbochsmevexperimentálnomsúbore(n=50)nameralivstupných 75,31 ± 8,20 stupna a výstupných 78,42 ± 8,35 stupna. U probandov kontrolnej skupiny boltentoposunz76,51±8,69na79,28±8,43stupna.VýsledkyporovnávaGraf12. 100,00 80,00 75,31 76,51 78,42 79,28 60,00 40,00 20,00 0,00 Predterapiou Poterapii )enputs(ubyhophaszoR GoniometriaPIPklbov--Flexia ITF+kinezioterapia Kinezioterapia Graf12GoniometriaPIP--Flexia(stupne)uobochskupín 41 Tabul’ka6Goniometriaklbovruky--MCPaPIP(stupne) Parameter/Strana x¯Pred sdPred min-max x¯Po sdPo min-max MCP-Flexia(ExperimentálnaskupinaITF) Dx(vpravo) 67,82 7,24 55-82 73,06 7,32 61-90 Sin(vl’avo) 66,78 6,84 52-80 72,14 7,28 58-88 MCP-Flexia(KontrolnáskupinaK) Dx(vpravo) 68,62 8,31 56-84 70,42 7,81 58-83 Sin(vl’avo) 67,18 8,83 53-86 69,13 8,58 55-86 MCP-Extenzia(ExperimentálnaskupinaITF) Dx(vpravo) 35,08 6,28 22-54 39,26 6,15 25-58 Sin(vl’avo) 34,76 6,82 20-52 39,16 6,47 23-56 MCP-Extenzia(KontrolnáskupinaK) Dx(vpravo) 34,32 6,83 23-55 35,52 6,78 25-54 Sin(vl’avo) 34,14 6,82 23-54 35,40 7,08 22-54 PIP-Flexia(ExperimentálnaskupinaITF) Dx(vpravo) 74,50 8,12 62-95 77,98 8,45 65-102 Sin(vl’avo) 76,12 8,28 64-98 78,86 8,25 68-105 PIP-Flexia(KontrolnáskupinaK) Dx(vpravo) 76,90 8,72 64-92 79,66 8,56 67-95 Sin(vl’avo) 76,12 8,66 64-92 78,89 8,29 68-95


\subsection{Analýzagoniometrickýchmeraníklbovnohy(MTP)}

Zásadným prínosom nášho výskumu je sledovanie pohyblivosti MTP klbov (meta- tarzofalangeálnych),ktorésúpriRAcastopostihnuté"prednosrdímädeformitami. V experimentálnej skupine (n=50) sme pri flexii v MTP namerali pred terapiou priemernú hodnotu 31,42 ± 4,25 stupna. Po liecbe sa rozsah zvýšil na 39,82 ± 4,15 stupna. U kontrolnej skupiny (n=50) bol progres pri flexii v MTP z 31,09 ± 6,42 na 32,19 ± 6,05 stupna.VýsledkydokumentujeGraf13. 42 100,00 80,00 60,00 40,00 39,82 31,42 31,09 32,19 20,00 0,00 Predterapiou Poterapii )enputs(ubyhophaszoR GoniometriaMTPklbov--Flexia ITF+kinezioterapia Kinezioterapia Graf13GoniometriaMTP--Flexia(stupne)uobochskupín Pri extenzii v MTP v experimentálnej skupine (n=50) narástol rozsah pohybu zvstupných55,70±6,30stupnanavýstupných66,14±6,10stupna.Vkontrolnejskupine (n=50) sme namerali zmenu z 58,41 ± 5,86 na 60,41 ± 5,86 stupna. Tieto výsledky (Graf 14)svedciaovýznamnombenefitekombinovanejterapieajprestabilitustojaachôdze. 100,00 80,00 66,14 60,00 60,41 58,41 55,70 40,00 20,00 0,00 Predterapiou Poterapii )enputs(ubyhophaszoR GoniometriaMTPklbov--Extenzia ITF+kinezioterapia Kinezioterapia Graf14GoniometriaMTP--Extenzia(stupne)uobochskupín 43 Tabul’ka7Goniometriaklbovnohy--MTP(stupne) Parameter/Strana x¯Pred sdPred min-max x¯Po sdPo min-max MTP-Flexia(ExperimentálnaskupinaITF) Dx(vpravo) 31,48 4,32 22-38 39,85 4,20 30-48 Sin(vl’avo) 31,36 4,18 22-40 39,79 4,10 31-48 MTP-Flexia(KontrolnáskupinaK) Dx(vpravo) 31,40 6,84 22-42 32,46 6,10 25-42 Sin(vl’avo) 30,78 6,00 22-42 31,92 6,01 22-43 MTP-Extenzia(ExperimentálnaskupinaITF) Dx(vpravo) 55,75 6,38 45-66 66,18 6,15 55-78 Sin(vl’avo) 55,65 6,22 44-68 66,10 6,05 54-78 MTP-Extenzia(KontrolnáskupinaK) Dx(vpravo) 58,26 6,36 45-65 60,18 6,28 45-68 Sin(vl’avo) 58,56 5,37 48-65 60,64 5,45 50-68 Komplexnýpohl’adnasubjektívnystavpacientovposkytolpodrobnýdotazníkvlast- nej konštrukcie, ktorý sledoval 10 kl’úcových parametrov (ADL, bolest’ pri chôdzi, stu- hnutost’, únava, bolesti v klboch, mobilita, motorika, úzkost’, spánok a práca). V experi- mentálnej skupine (ITF) bolo priemerné vstupné skóre 59,20 ± 18,56 bodu a po liecbe kleslona49,22±18,52bodu.Vkontrolnejskupinesmezaznamenalimiernejšiezlepšenie z50,30±4,19na46,06±5,23bodu.PorovnanieobochskupínznázornujeGraf15. 80,00 60,00 59,20 49,22 50,30 46,06 40,00 20,00 0,00 ITF+kinezioterapia Kinezioterapia vodobtecúS Celkovéskóredotazníkavlastnejkonštrukcie Predterapiou Poterapii Graf15Celkovéskóredotazníkavlastnejkonštrukcieuobochskupín 44


\subsection{Verifikáciastanovenýchhypotéz}

Nazákladepodrobnéhoštatistickéhospracovaniaarozpísaniadátuvedenýchvsekcii 3.6(Tabul’ky3--7)smepristúpilikvyhodnoteniustanovenýchhypotéz. Hypotéza c. 1: Predpokladáme, že pacienti s reumatoidnou artritídou, ktorí ab- solvujú terapiu interferencnými prúdmi v kombinácii s kinezioterapiou, uvádzajú niž- šiuintenzitubolestinežpacientiliecenílenkinezioterapiou. Na overenie tejto hypotézy sme použili výsledky štatistického spracovania intenzity bolesti zaznamenanej v dotazníku (Tabul’ka 3). V experimentálnej skupine (n=50) bola priemernáhodnotabolestiprizát’ažipredterapiou7,24±2,10bodovapoukonceníkombi- novanejliecbysmedosiahlivýraznézlepšeniena4,92±1,72bodov.Vkontrolnejskupine (n=50) klesla priemerná hodnota bolesti pri zát’aži z vstupných 7,00 ± 2,06 na výstupných 5,40 ± 1,76 bodov. Pri porovnaní oboch skupín vidíme, že v skupine s elektroliecbou (ITF) došlo k poklesu o 2,32 bodu, zatial’ co v kontrolnej skupine len o 1,6 bodu. Vyššie uvedená analýza potvrdzuje väcší analgetický efekt kombinovanej terapie. Na základe tohto hodno- teniahypotézuH1prijímame. Hypotézac.2:Predpokladáme,žekombináciainterferencnýchprúdovakinezi- oterapiemáväcšípozitívnyvplyvnarozsahpohybuvpostihnutýchklbochakosamotná kinezioterapia. Kverifikáciidruhejhypotézysmevyužiliúdajezgoniometrickýchmeraní(Tabul’ky 5,6a7).Prisledovaníflexievzápästíuexperimentálnejskupinyvzrástolrozsahz60,14±6,94 na65,50±6,96stupna,kýmvkontrolnejskupinebolposunz59,54±7,05na61,49±7,14 stupna. Pri extenzii v zápästí sme zaznamenali v ITF skupine progres z 50,68 ± 7,46 na 55,79±7,47stupna,zatial’cokontrolnáskupinavykazovalamiernejšiezlepšeniez52,18±7,13 na54,11±7,16stupna. Obdobnýtrendsmezaznamenaliajpridrobnýchklbochruky.FlexiavMCPklboch v experimentálnom súbore vykázala zlepšenie o 5,30 stupna, zatial’ co v kontrolnom o 1,88 stupna. Pri extenzii v MCP klboch bol posun u ITF skupiny 4,29 stupna oproti 1,23 stupna v kontrolnom súbore. Flexia v PIP klboch vykazovala u experimentálnej skupiny nárast z75,31±8,20na78,42±8,35stupnaaukontrolnejskupinyz76,51±8,69na79,28±8,43 stupna. Najmarkantnejšie rozdiely sme pozorovali pri klboch nohy. Flexia v MTP klboch uITFskupinydosiahlaprogreso8,40stupna,kýmukontrolnejleno1,10stupna.Priexten- zii v MTP klboch sme u experimentálnej skupiny namerali nárast o 10,44 stupna (z 55,70 na 66,14), zatial’ co u kontrolnej skupiny bol tento prírastok len 2,00 stupna. Všetky sle- dované segmenty vykazujú u skupiny s interferencnými prúdmi výraznejší nárast klbovej pohyblivosti.NazákladetýchtozisteníhypotézuH2prijímame. Hypotéza c. 3: Predpokladáme, že pacienti liecení kombinovanou terapiou sub- jektívne hodnotia kvalitu svojho života pozitívnejšie ako pacienti v skupine bez elek- troterapie. 45 Na overenie tretej hypotézy sme analyzovali parametre rannej stuhnutosti (Tabul’ka 4) a celkové skóre dotazníka vlastnej konštrukcie (Graf 15). Intenzita rannej stuhnutosti vexperimentálnejskupinekleslaz6,02±2,24na3,02±1,95bodu(zlepšenieo3body),pri- comvkontrolnejskupinebolpoklesonecelý1bod(5,32na4,42).Celkovéskóredotazníka (kde nižšia hodnota znamená lepšiu kvalitu života) kleslo u ITF skupiny z 59,20 ± 18,56 na 49,22 ± 18,52 bodu. U kontrolnej skupiny bol tento pokles ovel’a miernejší (z 50,30 na 46,06).Výsledkypotvrdzujú,žepacientiskombinovanouliecbousubjektívnepocit’ujúmen- šiefunkcnéobmedzeniealepšiezvládajúbežnédennéaktivity.Nazákladetýchtovýsledkov hypotézuH3prijímame.


\section{Etickéaspektyvýskumu}

Výskumprebiehalvsúladesetickýmiprincípmiklinickéhoskúšania.Všetcipartici- pujúcipacientibolivopredoboznámenísciel’omprieskumuadobrovol’nesúhlasilisposkyt- nutím údajov pre vedecké úcely. Zber dát bol realizovaný anonymne, bez zaznamenávania osobných identifikacných údajov, cím bola zabezpecená plná ochrana súkromia všetkých probandov.Získanévýsledkyslúžiavýhradnenaúcelyspracovaniatejtodiplomovejpráce. 46 4 Diskusia V súcasnej dobe predstavuje reumatoidná artritída (RA) jednu z najzávažnejších vý- ziev modernej reumatológie a rehabilitácie. Ide o systémové zápalové ochorenie, ktoré svo- jou progresiou vedienielen k štrukturálnym zmenám naklboch, ale výraznedeštruuje aj so- ciálneväzbyapracovnúschopnost’pacienta.Socio-ekonomickázát’ažspojenásdlhodobou liecbou a prípadnou invalidizáciou je enormná. V našej diplomovej práci sme sa preto za- merali na hl’adanie optimálnych terapeutických ciest, ktoré by dokázali efektívne zmiernit’ symptómy a spomalit’ funkcný úpadok pacienta. Porovnanie kombinovanej terapie (inter- ferencné prúdy a kinezioterapia) oproti izolovanej kinezioterapii nám poskytlo cenné dáta, ktorépotvrdzujúpotrebumultidisciplinárnehoprístupu.


\section{Interpretáciavýsledkovabiofyzikálnesúvislosti}

Interpretáciou výsledkov sme dospeli k jednoznacnému potvrdeniu stanovených hy- potéz. Z hl’adiska biofyziky je kl’úcové pochopit’, preco práve interferencné prúdy (IFP) v kombinácii s pohybom vykazujú taký výrazný synergický efekt. IFP pôsobia v hlbke tka- niva,kdedochádzakinterferenciidvochstrednefrekvencnýchprúdov,címvznikánízkofrek- vencný modulacný úcinok. Tento mechanizmus umožnuje prekonat’ kožný odpor bez dráž- denia nociceptorov na povrchu, co je u reumatikov s precitlivenou kožou (casto vplyvom kortikoterapie)kl’úcovávýhoda. Vplyvom kombinovanej terapie na intenzitu bolesti (H1) sme zaznamenali v na- šomexperimentálnomsúborepoklespriemernejhodnotyVASprizát’ažio2,32bodu,zatial’ co v kontrolnej skupine bol tento pokles len 1,60 bodu. Tento rozdiel pripisujeme schop- nosti interferencných prúdov blokovat’ prenos bolestivých vzruchov cestou teórie vrátkovej kontroly(GateControlTheory),ktorúdefinovaliMelzackaWall.Podl’atejtoteóriestimulá- ciahrubýchaferentnýchvlákienA-beta„zatvárabránu“prebolestivéstimulyprichádzajúce cez tenké vlákna C. Navyše, aplikácia IFP vyvoláva vazodilatáciu a zlepšuje mikrocirkulá- ciu, co vedie k rýchlejšiemu odplavovaniu algogénnych látok (laktát, serotonín, histamín) zo zápalového ložiska. Fuentes et al. (2010) vo svojej meta-analýze uvádzajú, že práve táto „metabolická ocista“ tkaniva je zodpovedná za dlhodobejší analgetický efekt v porovnaní sbežnýmTENSprúdom. Pri hodnotení rozsahu pohybu (H2) boli výsledky v prospech experimentálnej sku- pinyeštemarkantnejšie.NajvýraznejšiezlepšeniesmepozorovalipriextenziivMTPklboch (nárast o 10,44 stupna). Tento jav vysvetl’ujeme tým, že potlacenie bolesti pomocou IFP umožnilo pacientom prekonat’ tzv. antalgický svalový spazmus. Reumatik podvedome ob- medzuje pohyb v klbe, aby sa vyhol bolesti, co vedie k sekundárnemu skracovaniu mäk- kých tkanív a klbových puzdier. Interferencné prúdy svojím jemným vibracným úcinkom pôsobia ako „mikromasáž“, cím uvol’nujú hypertonus svalov a zlepšujú viskoelasticitu ko- 47 lagénnych vlákien. Tým sa vytvára ideálne „okno príležitosti“ pre následnú kinezioterapiu. Naše výsledky korešpondujú s teóriou vrátkovej kontroly bolesti, ktorú vo svojom prehl’ade rozoberajú Érika Patrícia Rampazo et al. (2022). Zmiernenie bolesti umožnuje pacientovi vykonávat’ cvicenia vo väcšom rozsahu a s menším svalovým odporom. Cvicenie v od- bolestnenom stave je nielen efektívnejšie, ale pre pacienta aj psychicky prijatel’nejšie, co zvyšuje compliance k liecbe. Kolár (2009) vo svojej ucebnici zdôraznuje, že elektroterapia bez následného cvicenia je v rehabilitácii RA len polovicatým riešením, s cím sa na základe našichdátplnestotožnujeme. Subjektívne vnímanie kvality života (H3) sme analyzovali nielen cez rannú stuhnu- tost’, ale aj cez 10 špecifických ADL parametrov. Zlepšenie celkového skóre v ITF skupine (pokles o cca 10 bodov) jasne reflektuje lepšiu adaptáciu pacienta na bežný život. Zaují- mavým zistením bol vplyv na spánok a úzkost’. Pacienti uvádzali, že vd’aka ústupu nocnej bolesti sa zlepšila kvalita ich oddychu, co následne znižovalo dennú únavu -- jeden z hlav- nýchlimitujúcichfaktorovuRA.Rannástuhnutost’,ktorúpacientisRAvnímajúakovel’mi obmedzujúcu, klesla v experimentálnom súbore v pomere 3:1 oproti kontrole. Tento efekt pripisujemeresorpciiperikapsulárnychedémov,ktorúIFPstimulujúcestouzlepšenejlymfa- tickejdrenáže.


\section{Porovnaniesliterárnymipramenmi}

Náš výskum korešponduje s viacerými recentnými štúdiami. Zhang et al. (2025) vo svojej vedeckej práci zdôraznuje, že kombinácia fyzikálnych modalít a cvicenia je v liecbe reumatických ochorení efektívnejšia než izolované prístupy. Zhangov výskum sa síce pri- márne zameriaval na vel’ké klby, avšak princíp akcelerácie liecebného procesu pri použití synergiejeidentickýsnašimizisteniaminadrobnýchklbochrukyanohy. Podobne Alayat et al. (2016) potvrdili, že pridanie prístrojovej fyzikálnej terapie k liecebnej telesnej výchove (LTV) výrazne znižuje chronickú bolest’. Hoci vo svojej štúdii používali vysokovýkonný laser, mechanizmus tlmenia bolesti a podpory regenerácie tkanív je v mnohom analogický pôsobeniu interferencných prúdov. Náš výber IFP v praxi TnUAD bol motivovaný lepšou dostupnost’ou tejto modality a jej schopnost’ou plošne zasiahnut’ viacerédrobnéklbysúcasne(napr.metakarpy),cojeprilaserovejterapiicasovonárocnejšie. V kontraste s niektorými staršími prácami, ktoré spochybnovali význam elektrotera- pie u reumatikov, sa stotožnujeme s modernými názormi Metsios et al. (2020). Tí uvádzajú, že moderné prístupy k LTV musia byt’ podporené adekvátnou kontrolou bolesti, aby ne- došlo k nežiaducemu pret’ažovaniu zapálených štruktúr. Naše dáta z dotazníkov (Graf 15) potvrdzujú, že pacienti vnímajú kombinovanú terapiu ako bezpecnú a vysoko efektívnu, co jekl’úcovépredlhodobéudržanieliecebnéhorežimu. Pri hlbšom štúdiu domácej literatúry, konkrétne prác K. Pavelka (2012) a Karel Pa- velka et al. (2018), nachádzame potvrdenie dôležitosti vcasnej intervencie. Pavelka zdôraz- nuje, že „vcasné okno“ liecby nie je dôležité len pri farmakoterapii biologikami, ale rov- 48 naký význam má aj v rehabilitácii. My doplname, že práve pridanie interferencnej terapie umožnuje pacientom zacat’ s aktívnym pohybom skôr a v intenzívnejšej forme, než by im dovol’ovalasamotnárannábolest’.


\section{Bio-psycho-sociálnyaspektaklinickývýznam}

Nemožno opomenút’ ani psychologický aspekt terapie. Pacienti podstupujúci kom- binovanú liecbu mali pravidelný kontakt s fyzioterapeutom pocas dvoch fáz (elektroterapia aj cvicenie), co mohlo pôsobit’ ako motivacný faktor. Vedomie, že využívajú „moderný prí- strojovýprístup“,zvyšovalodôveruvterapiu(placeboefektsmenesledovali,alejehopodiel nemožno úplne vylúcit’). Fyzioterapeut tu nepôsobí len ako inštruktor cvicenia, ale aj ako edukátor,ktorýpacientaucírozumiet’svojejbolestiareagovat’najejzmeny. Zaujímavýmvedl’ajšímzistenímnašejpráce(otázkac.1vdotazníku)bolfakt,žetyp pracovného zat’aženia (fyzické vs. psychické) nemal štatisticky významný vplyv na mieru zlepšenia pod vplyvom ITF. Toto zistenie je povzbudivé pre klinickú prax, pretože nazna- cuje,žekombinovanáterapiajeuniverzálneúcinnápreširokéspektrumpacientovsRA,bez ohl’adu na ich profesijný profil. To vyvracia mýtus, že fyzikálna terapia je vhodná len pre starších,nepracujúcichpacientov.


\section{Limityasilnéstránkyštúdie}

Každý výskum má svoje limity a naša práca nie je výnimkou. Za hlavný limit pova- žujeme obmedzený casový horizont sledovania. Reumatoidná artritída je celoživotné ocho- renie a hoci sme zaznamenali signifikantné zlepšenie bezprostredne po 10-dnovom cykle, dlhodobýefekt(follow-uppo3alebo6mesiacoch)nebolpredmetomtejtopráce.Taktiežva- riabilita v užívaných antireumatikách (NSAID, biologická liecba) u jednotlivých pacientov mohla ciastocne maskovat’ alebo naopak umocnovat’ výsledky, hoci sme sa snažili o rovno- mernérozdeleniesúboru. Zasilnéstránkynašejštúdievšakpovažujeme: • Vysoká precíznost’ merania: Goniometrické meranie všetkých drobných klbov (Dx aj Sin) poskytlo omnoho detailnejší obraz o stave pacienta než bežné meranie len naj- viacpostihnutéhoklbu. • Homogenita súboru: Rozdelenie 100 pacientov (50 vs 50) s reumatoidným ocho- rením predstavuje robustný vzorec pre diplomovú prácu, co zvyšuje štatistickú silu výpovede. • Kombinácia objektívnych a subjektívnych metód: Prepojenie goniometrie s kom- plexným10-parametrovýmdotazníkomumožnujeholistickýpohl’adnapacienta. Záveromdiskusiemôžemekonštatovat’,žepredloženévýsledkytvoriapevnýzáklad pre odporúcanie kombinovanej terapie do bežnej rehabilitacnej praxe. Synergia medzi hlb- kovým pôsobením interferencných prúdov a cielenou kinezioterapiou sa ukazuje ako kl’úc 49

